\section*{Introduction}
{Deduced} in 1871, the Saint-Venant equations \cite{SaintVenant} (or 1-D shallow water equations) are among the most famous equations in fluid dynamics and have been investigated in hundreds of studies.
Although being quite simple, their richness has made them become a major tool in practice for many industrial goals, the most famous being probably the regulation of navigable rivers.
They are the ground model for such purpose in France and Belgium.
Regulation of rivers
is a major issue, for navigation, freight transport,
renewable energy production,
but also for safety reasons, especially as
several nuclear plants all around the world are implanted close to rivers.
For these reasons, the stability of the steady-states of the Saint-Venant equations has been, and is still, a major issue.

Many results were obtained in the last decades.
In 1999, the robust stability of the homogeneous linearized Saint-Venant equations was shown using a Lyapunov approach and proportional feedback controllers \cite{coron1999}.
Later, the stability of the homogeneous nonlinear Saint-Venant equations was achieved, still using proportional feedback controllers.
In 2008, through a semi-group approach \cite{dos2008}, the stability of the inhomogeneous nonlinear Saint-Venant equation was shown for sufficiently small friction and slope (or equivalently sufficiently small canal), and these results were successfully applied to real data sets from the Sambre river in Belgium.
More recently, in~\cite{BC2017} the authors have given sufficient conditions to stabilize the nonlinear Saint-Venant equations with arbitrary friction for the $H^{2}$ norm but no slope using again proportional feedback controllers,
and in \cite{HS} with both arbitrary friction and slope.
This last result is proved by exhibiting {a Lyapunov function that has a simple form close to a local entropy} for the nonlinear inhomogeneous Saint-Venant equations.

It is worth mentioning that other stability results have also been obtained in less classical cases or with less classical feedback laws.
For instance, in \cite{stvenantjump} was shown the rapid stabilization of the homogeneous nonlinear Saint-Venant equations when a shock (\eg a hydraulic jump) occurs in the target steady-state.
Such {a} shock indu\-ces new difficulties and the presence of shocks can limit in general the controllability and the stability in weaker norms of hyperbolic systems
with boundary controls
\hbox{\cite{ABG,BG}}.
Also, several results (e.g.
\cite{CVBK})
were obtained using a backstepping approach, a very powerful method based on a Volterra transformation, developed mainly for PDE in~\cite{KrsticBook},
and generalized recently with a Fredholm transformation for hyperbolic systems \hbox{\cite{CoronFredholm,Zhang,Zhang2}}.
One may look at \cite{HS} for a more detailed survey about this method and its use for the Saint-Venant equations.
However, backstepping gives rise to non-local and non-static feedback laws that are likely to be harder to implement, and, to our knowledge, have not been implemented yet.

Most of the previous results were performed with static proportional feedback controllers. When it comes to industrial applications, however, the proportional integral~(PI) control is by far the most popular regulator.
It is used for instance for the regulation of the Sambre and Meuse river in Belgium \cite[Chap.\,8]{BastinCoron1D}.
The reason behind such preference is the robustness of the PI control with off-set errors \cite[Chap.\,11.3]{Astrom2}. {An example can be found in \cite{DosSantosBastinCoron}
where the authors show the interest of adding an integral term to a proportional control on a linear and homogeneous system, and exhibit coherent experimental result.}\enlargethispage{.5\baselineskip}%

For these reasons, the PI controller has fed a wide literature, at least when used on finite dimensional systems.
However, despite their indisputable practical interest, PI controllers for nonlinear infinite dimensional systems have shown hard to handle mathematically
and even studying simple systems give sometimes rise to lengthy proofs with relatively sophisticated tools \cite{PItransport}.
While the behaviour and the stability of linearized equations
with PI controller
has been well understood in the past, partly thanks to spectral tools like the spectral mapping theorem (\eg \cite{Neves,Lichtner} for hyperbolic systems), no such {tool exists} for nonlinear systems, and
the stability of the nonlinear Saint-Venant equations has remained a challenge until today.
{Among the existing linear results using a spectral approach, one can refer to \cite{XuSallet,XuSallet2014} where the authors find a sufficient condition for the stabilization of the linearized inhomogeneous Saint-Venant equations.
Necessary and sufficient conditions for the linearized homogeneous Saint-Venant equations
are given in \cite[\S\S 2.2.4.1,\,3.4.4]{BastinCoron1D}.
In \cite{CoronTamasoiu2015} the authors find a necessary and sufficient condition for a linear scalar equation and show the difficulty of finding good conditions for the nonlinear equation, while
in \cite{BastinCoronTamasoiu2015} the authors deal with $2\times2$ systems.}
Among the existing nonlinear results
{one can refer to {\cite{TerrandAndrieuDos} in the case where the operator without PI control generates an exponentially stable semi-group},
\cite{Trinh2015} where the authors find a sufficient condition for the nonlinear homogeneous Saint-Venant equations,
\cite[2.2.4.2]{BastinCoron1D} where the authors find a necessary and sufficient condition also for the
nonlinear homogeneous Saint-Venant equations,
while
\cite[\S\S 5.4.4,\,5.5]{BastinCoron1D} and \cite{BastinCoron2013} give a sufficient condition for the inhomogeneous Saint-Venant equations for a single channel or a network,
but in the particular case of constant steady-states only, which simplifies their analysis \cite{C1}. Strictly speaking, this last result was derived for the linearized system
but with a Lyapunov approach, which can easily be generalized to the nonlinear system.}
More recently, and this is the most advanced result yet, \cite{BastinCoronPI} gave a sufficient condition of stability for the inhomogeneous Saint-Venant equations with an arbitrary friction and river length but only
in the absence of slope, using a Lyapunov approach.

In this paper, we consider the stabilization of the general nonlinear Saint-Venant equations with a single boundary PI control.
We give a simple and explicit condition on the parameters of the PI controller such that any steady-state
is exponentially stable for the $H^{2}$ norm.
While stability results in inhomogeneous and nonlinear systems often imply a limit length for the domain, depending on the source term, above with we are unable to guarantee any stability (\cite{C1,C1_22,BastinCoron22,dos2008} or \cite[Chap.\,6]{BastinCoron1D}),
this result holds whatever the friction, the slope, and the length of the channel.
Besides, our condition is independent of the slope, the friction coefficient, the river length, and, more surprisingly, {can be made independent of }the steady-state considered.
Finally, when there is no slope this condition is less restrictive than the condition obtained in~\cite{BastinCoronPI}
and when there is no friction or slope this condition coincides with the {necessary and sufficient} spectral condition of stability for the linearized system given in \cite{BastinCoronTamasoiu2015} and \cite[Th.\,2.7]{BastinCoron1D}.\enlargethispage{.5\baselineskip}

The case where the inflow disturbances are time dependent and no steady-states exists was seldom considered in the literature.
However, it is in fact unlikely that the industrial target state is a real steady-state as the inflow disturbance often depends on time in practice,
even though only slowly.
Therefore,
in the more general framework of slowly time-varying target states,
we show the Input-to-State Stability (ISS) of the system with respect to the variation of the inflow disturbance.
Finally, we show that if
we allow the controller to depend on the target state,
by changing slightly the PI controller,
we can ensure
the exponential stability of slowly-varying target trajectories. {These trajectories} are the natural {targets} to consider when there is no steady-state of the system.

This paper is organized as follows: in Section \ref{s1} we give a description of the nonlinear Saint-Venant equations, we introduce the time-varying target trajectories together with some definitions and existence results, then we state our main results.
In Section~\ref{s2} we prove our main result, Theorem \ref{th1}, that deals with the exponential stability of time-varying state.
In the appendix, we show that Corollary \ref{cor1} dealing with the exponential stability of steady-states,
and Theorem \ref{th2} showing the ISS of the system with respect to the variation of the inflow disturbance,
are both deduced from the proof of Theorem \ref{th1}.

\subsubsection*{Acknowledgments}
The author would like to thank Jean-Michel Coron for his constant support and his advices, and Sebastien Boyaval for many fruitful discussions.
The author also wishes to thank Eric Demay, Peipei Shang, Shengquan Xiang and Christophe Zhang for fruitful discussions.

\section{Model description}
\label{s1}
We
consider
the following nonlinear Saint-Venant equations for a rectangular channel with arbitrary slope and friction.\vspace*{-3pt}
\begin{equation}
\begin{split}
\partial_{t}H+\partial_{x}(HV)&=0,\\[-3pt]
\partial_{t}V+V\partial_{x}V+g\partial_{x}H+\Bigl(\frac{kV^{2}}{H}-C(x)\Bigr)&=0.
\end{split}
\label{Stvenant}
\end{equation}
Here, $k$ is an arbitrary nonnegative friction coefficient and $C$ denotes the slope, which is assumed to be a $C^{2}$ function, with $C(x):=-gdB/dx$ where $B$ is the bathymetry and $g$ the acceleration of gravity.
We are interested in systems where the water flow uphill is a given function, unknown and imposed by external conditions, for instance a flow coming from another country, while the water flow downhill is controlled through a hydraulic installation. Therefore, we have the following boundary conditions,
\begin{equation}
\begin{split}
H(t,0)V(t,0)&=Q_{0}(t),\\
H(t,L)V(t,L)&=U(t),
\end{split}
\label{boundary02}
\end{equation}
where $U(t)$
is a control feedback and $Q_{0}(t)$
is the incoming flow, which is a given (and unknown) function. Here $L$ denotes the length of the water channel.
In practical situations, the formal control $U(t)$ can be expressed by a simple linear model \cite{BastinCoronPI}
\begin{equation}
U(t)=v_{G}(H(t,L)-U_{1}(t)),
\label{vG}
\end{equation}
where $U_{1}(t)$ is the elevation of the gate of the dam, which is the real control input that can be chosen, while $v_{G}$ is a constant depending on the parameters of the gate (potentially unknown as well).

\subsection{Control goal and target trajectory}
Usually, the industrial goal of such system is to stabilize the level of the water at the end point $H(t,L)$, called control point, to a target value $H_{c}>0$.
On the other hand, the usual mathematical goal in such a problem is to stabilize a target steady-state $(H^{*},V^{*})$, potentially nonuniform \cite[Preface]{BastinCoron1D}. However, in the present problem \eqref{Stvenant}--\eqref{boundary02}, it is clear that, when $Q_{0}$
is not constant, it is impossible to aim at stabilizing any steady-state and one needs to aim at stabilizing other target trajectories.
Therefore, we define the following target trajectory $(H_{1},V_{1})$ that we aim to stabilize as the solution of
\begin{equation}\label{target}
\begin{split}
\partial_{t}H_{1}+\partial_{x}(H_{1}V_{1})&=0,\\
\partial_{t}V_{1}+V_{1}\partial_{x}V_{1}+g\partial_{x}H_{1}+\Bigl(\frac{kV_{1}^{2}}{H_{1}}-C(x)\Bigr)&=0,\\
H_{1}(t,0)V_{1}(t,0)&=Q_{0}(t),\\
H_{1}(t,L)&=H_{c},
\end{split}
\end{equation}
with the initial condition
\begin{equation}\label{initialtarget}
H_{1}(0,\cdot)=H^{*}(\cdot)\quad\text{and}\quad V_{1}(0,\cdot)=V^{*}(\cdot),
\end{equation}
where $(H^{*},V^{*})$ is the (unique) steady-state solution of the system {when $Q_{0}$ is} constant, equal to $Q_{0}(0)$. Namely, $(H^{*},V^{*})$ is the solution of
\begin{equation}
\label{targetsteady0}
\begin{split}
\partial_{x}(HV)&=0,\\
V\partial_{x}V+g\partial_{x}H+\Bigl(\frac{kV^{2}}{H}-C(x)\Bigr)&=0,\\
H(L)&=H_{c},
\end{split}
\end{equation}
with condition at $x=0$
\begin{equation}
\begin{split}
&H^{*}(0)V^{*}(0)=Q_{0}(0).
\end{split}
\label{targetsteady}
\end{equation}
We are now going to show that the trajectory $(H_{1},V_{1})$ exists for any time and satisfies some bounds.

\subsubsection*{Existence and bounds of the target trajectory $(H_{1},V_{1})$}

Instead of studying directly our target trajectory $(H_{1},V_{1})$ we first construct an intermediary family of functions
$(H_{0},V_{0})$ where at each time $t$, $(H_{0}(t,\cdot),V_{0}(t,\cdot))$ is defined as the space dependent steady-state that would be associated with the constant flow $Q_{0}(t)$. This is detailed in the following paragraph.

We defined previously $(H^{*}, V^{*})$ as the steady-state associated to a constant flux $Q_{0}\equiv Q_{0}(0)$, {that is} $(H^{*},V^{*})$ is the solution of the ODE problem \eqref{targetsteady0} with
initial condition $H^{*}(0)V^{*}(0)=Q_{0}(0)$. But in fact
at each time $t^{*}\in\mathbb{R}^{*}_{+}$, we can also define a steady-state $(H^{*}_{t^{*}},V^{*}_{t^{*}})$ associated to a constant flux $Q_{0}\equiv Q_{0}(t^{*})$. {In other words}
$(H^{*}_{t^{*}},V^{*}_{t^{*}})$ is the solution of the ODE problem \eqref{targetsteady0}
with initial condition satisfying
\begin{equation}
\begin{split}
&H_{t^{*}}^{*}(0)V_{t^{*}}^{*}(0)=Q_{0}(t^{*}).
\end{split}
\label{targetsteady2}
\end{equation}
Although the system \eqref{targetsteady0}, \eqref{targetsteady2} could seem peculiar as it has boundary conditions imposed both in $0$ and in $L$,
we know looking at the first equation of \eqref{targetsteady0} that this system \eqref{targetsteady0}, \eqref{targetsteady2} is in fact equivalent to a single ODE on $H_{t^{*}}^{*}$
with boundary condition $H_{t^{*}}^{*}(L)=H_{c}$ and the function $V_{t^{*}}^{*}$ defined by $V_{t^{*}}^{*}=Q_{0}(t^{*})/H_{t^{*}}^{*}$. Indeed the first equation of \eqref{targetsteady0} is equivalent to saying that $H_{t^{*}}^{*}V_{t^{*}}^{*}$ is a constant function, equal to $Q_{0}(t^{*})$ thanks to \eqref{targetsteady2}. In other words \eqref{targetsteady0}, \eqref{targetsteady2} is equivalent to
\begin{equation*}
\begin{cases}
V_{t^{*}}^{*}(x)=\dfrac{Q_{0}(t^{*})}{H_{t^{*}}^{*}(x)},\quad\forall x\in[0,L],\\[8pt]
\Bigl(g-\dfrac{Q_{0}^{2}(t^{*})}{H_{t^{*}}^{*3}} \Bigr)\partial_{x}H_{t^{*}}^{*}+ \Bigl(\dfrac{k^{2}Q_{0}(t^{*})^{2}}{H_{t^{*}}^{*3}}-C(x) \Bigr)=0,\\[5pt]
H(L)=H_{c}.
\end{cases}
\end{equation*}
Thus for each $t^{*}\in[0,+\infty)$ such function exists on $[0,L]$, is unique and $C^{3}$ provided that the state stays in the fluvial regime (or subcritical regime), \ie $g H_{t^{*}}^{*}>V^{*2}_{t^{*}}$ on $[0,L]$.
This, for a given $H_{c}$, is equivalent to a bound on $Q_{0}(t^{*})$
(see \cite{HS} for more details).
As we are interested in {stabilizing} physical trajectories in the fluvial regime, we assume that this assumption is satisfied {in the following} and that
there exist $\alpha>0$ and $H_{\max}>0$ independent of $t^{*}\in[0,\infty)$ such that
\begin{equation}
\begin{split}
H_{t^{*}}^{*}<\frac{1}{2}H_{\max}\quad\text{on }[0,L],\\
gH_{t^{*}}^{*}-V_{{t^{*}}}^{*2}>2\alpha\quad\text{on }[0,L].
\end{split}
\label{fluvial0}
\end{equation}
For a given $H_{c}$, this is again equivalent to imposing a bound $Q_{\infty}$ on $\lVert Q_{0}\rVert_{L^{\infty}(0,\infty)}$, from \eqref{targetsteady0} and \eqref{targetsteady2},
(which is more logical {from an applicative point of view}).
However, for convenience, we will still use $H_{\max}$ and $\alpha$ in the following. This assumption is quite physical: in practical situation the river is in fluvial regime and $Q_{0}(t)$ is often periodic or quasi-periodic.
This gives a family of one-variable functions indexed by a~parameter $t^{*}$, which can also be seen as the two-variable functions
\begin{equation*}
(H_{0},V_{0}) : (t,x)\to (H^{*}_{t}(x),V^{*}_{t}(x)).
\end{equation*}
Besides, from \eqref{targetsteady}, as $(H_{t}^{*},V_{t}^{*})$
is the solution of a system of ODE with a parameter~$t$, the two variable functions $(H_{0},V_{0})$ therefore belongs to $C^{3}([0,+\infty)\times(0,L))$ (see \cite[Chap.\,5, Cor.\,4.1]{Hartman}).
{From} its definition, one can note that
\[
(H_{0}(0,\cdot),V_{0}(0,\cdot))=(H^{*},V^{*}).
\]

For clarity, we summarize here the different families of functions we introduced.
\begin{itemize}
\item
$(H^{*},V^{*})$, a function of $x$, the steady-state of the system when $Q_{0}\equiv \text{const.}$
\item
$(H_{1},V_{1})$, a function of $t$ and $x$, the target trajectory to reach when $Q_{0}$ is not a
constant. This trajectory is compatible with the objective $H(t,L) = H_{c}$, for any $t\in [0,T]$.
\item
$(H^{*}_{t^{*}},V^{*}_{t^{*}})$, a function of $x$, the steady-state of the system when $Q_{0}$ is a constant equal to $Q_{0}({t^{*}})$ ($t^{*}$ is fixed).
\item
$(H_{0},V_{0})$, a function of $t$ and $x$, the family such that
\[
(H_{0}, V_{0}):(t,x) \to (H_{t}^{*}(x),V^{*}_{t}(x)).
\]
\end{itemize}

Now that we have introduced this intermediary family of functions, we can show the existence of the target trajectory $(H_{1},V_{1})$ and
we have the following Input-to-State Stability (ISS) result {(see \cite{Sontag} for a definition of ISS for finite dimensional systems,
\cite[Chap.\,1, Chap.\,3]{IassonKrstic} for a generalization to first-order hyperbolic PDE and \cite{PrieurISS} for the use of Lyapunov function to achieve ISS on time-varying hyperbolic systems)},
\begin{prop}\label{propISS}
Assume that $\partial_{t}Q_{0}\in C^{2}([0,\infty))$. There exist positive constants $c_{1}$, $c_{2}$,
$\mu>0$, $\nu>0$ and $\delta>0$ such that if
$\lVert\partial_{t}Q_{0}\rVert_{C^{2}([0,+\infty))}\leq\delta$, then for any $(H_{1}^{0},V_{1}^{0})\in H^{2}((0,L),{\mathbb{R}^{2}})$ such that
\begin{equation*}
\lVert H_{1}^{0}-H^{*}\rVert_{H^{2}(0,L)}+\lVert V_{1}^{0}-V^{*}\rVert_{H^{2}(0,L)}\leq \nu,
\end{equation*}
the system \eqref{target} with initial condition $(H_{1}^{0},V_{1}^{0})$ has a unique solution $(H_{1},V_{1})\in C^{0}([0,+\infty),H^{2}(0,L))$ which satisfies the following ISS inequality
{\begin{multline}
\label{ISS}
\lVert H_{1}(t,\cdot)-H_{0}(t,\cdot)\rVert_{H^{2}(0,L)}+\lVert V_{1}(t,\cdot)-V_{0}(t,\cdot)\rVert_{H^{2}(0,L)}\\
\leq c_{1}(\lVert H_{1}^{0}-H^{*}\rVert_{H^{2}(0,L)}+\lVert V_{1}^{0}-V^{*}\rVert_{H^{2}(0,L)})e^{-\sfrac{\mu t}{2}}\\
+c_{2}\left(\int_{0}^{t}\left(|\partial_{t}Q_{0}(s)|+|\partial_{tt}^{2}Q_{0}(s)|+|\partial_{ttt}^{3}Q_{0}(s)|\right)e^{\sfrac{\mu s}{2}}ds\right)e^{-\sfrac{\mu t}{2}}.
\end{multline}}
\end{prop}
This result is shown in Appendix \ref{AppendixISS}, {and a definition of the $C^{2}$ norm is recalled in Remark \ref{rmkHpISS}}. Note that $Q_{0}$ is supposed to be bounded, which is quite physical, but there is no additional requirement on this bound besides the physical assumption given by $Q_{\infty}$ of remaining in the fluvial regime. This is important as in practical
situations the value of the incoming flow can change a lot, even though slowly.

Here, we {choose} to stabilize the trajectory $(H_{1},V_{1})$ associated to $H_{1}^{0}=H^{*}$ and $H_{1}^{0}=V^{*}$. As we will see, this target trajectory can be seen as the natural trajectory to stabilize as it satisfies the industrial goal $H(t,L)=H_{c}$
and it coincides with the steady-state solution when $Q_{0}$ is a constant.
{In this last case $Q_{0}$ and $H_{c}$ are imposed and $H^{*}$ and $V^{*}=Q_{0}/H^{*}$ are thus fully determined using \eqref{targetsteady0}.}
But one can note from \eqref{ISS} that, in fact, the behavior of $(H_{1},V_{1})$ at large time does not depend on the initial condition $(H_{1}^{0},V_{1}^{0})$ in \eqref{initialtarget}, provided that it is close in $H^{2}$ norm to $(H^{*},V^{*})$.

\begin{rmk}
\label{rmkHpISS}
The same ISS result can be shown replacing the $H^{2}$ norm in Proposition \ref{propISS} by the $H^{p}$ norm where ${p\in\mathbb{N}^{*}\setminus\{1\}}$, with the condition $\lVert\partial_{t}Q_{0}\rVert_{C^{p}([0,+\infty))}\leq\delta$ instead of $\lVert\partial_{t}Q_{0}\rVert_{C^{2}([0,+\infty))}\leq\delta$. This is shown in Appendix \ref{AppendixISS}. {We define here the $C^{p}$ norm for a function $U\in C^{p}(I)$, where $I$ is an interval, as
\begin{equation*}
\rVert U \lVert_{C^{p}(I)}:=\max_{i\in[0,p]}(\lVert \partial_{t}^{i}U\rVert_{L^{\infty}(I)})
\end{equation*}}
\end{rmk}

Thus, from Proposition \ref{propISS} and \eqref{fluvial0}, there exists a constant ${\delta}>0$ such that, if~$\lVert\partial_{t}Q_{0}\rVert_{{C^{2}([0,\infty))}}<{\delta}$, then $(H_{1},V_{1})\in C^{0}([0,+\infty),H^{2}(0,L))$ and
\begin{gather}
H_{1}(t,x)<H_{\max},\quad\forall(t,x)\in[0,+\infty)\times[0,L],
\label{Hinfty}\\
gH_{1}(t,x)-V_{1}^{2}(t,x)>\alpha,\quad\forall (t,x)\in[0,+\infty)\times[0,L].
\label{fluvial}
\end{gather}

Besides, when $Q_{0}$ is a constant, it is easy to check that $(H_{0},V_{0})=(H^{*},V^{*})$ is also solution of \eqref{target}--\eqref{initialtarget}. Thus, from the uniqueness of the solution of \eqref{target}--\eqref{initialtarget},
$(H_{1},V_{1})=(H^{*},V^{*})$ and, therefore, we recover a steady-state.
This illustrates that $(H_{1},V_{1})$ can be seen as the natural target state
when $Q_{0}$ is not a constant anymore.
Moreover, from \eqref{target}, stabilizing $(H_{1},V_{1})$ also satisfies the industrial goal by stabilizing $H(t,L)$ on the value $H_{c}$.

\subsection{Control design and main result}
As mentioned in the introduction, a usual type of controller used in practice to reach this aim is the proportional-integral (PI) controller.
It has the advantage of eliminating the offset coming from constant load disturbances, which can usually appear in these systems
as the command on the gate's level are only known up to some constant uncertainties.
A generic PI controller is given~by
\begin{equation}
U_{1}(t)=k_{p}(H_{c}-H(t,L))+k_{I}Z,
\label{defu1}
\end{equation}
where $k_{p}$ and $k_{I}$ are coefficients that can be designed and $Z$ accounts for the integral term, \ie
\begin{equation}
\dot Z=H_{c}-H(t,L).
\label{Z}
\end{equation}
With such controller, and using \eqref{vG}, the boundary conditions \eqref{boundary02} become {\eqref{Z} and}
\begin{equation}
\label{boundary2}
\begin{split}
H(t,0)V(t,0)&=Q_{0}(t),\\
H(t,L)V(t,L)&=v_{G}(1+k_{p})H(t,L)-v_{G}k_{p}H_{c}-v_{G}k_{I}Z.
\end{split}
\end{equation}
In Corollary \ref{cor1} we show that this boundary control can be used to stabilize {exponentially} a steady-state when $Q_{0}$ is a constant.
In Theorem \ref{th2} we show that this control can also provide an Input-to-State Stability property with respect to $\partial_{t}Q_{0}$.
However, this control \eqref{defu1} cannot be used to stabilize {a} dynamic target trajectory $(H_{1},V_{1})$, {whatever the coefficients,} as there is no function $Z_{1}\in C^{1}([0,+\infty))$ such that $(H_{1},V_{1},Z_{1})$
is a solution of \eqref{Stvenant}, \eqref{Z}, \eqref{boundary2} while $(H_{1},V_{1})$ is a solution of \eqref{target}. Therefore, when stabilizing a dynamic target trajectory, one has to add an additional term and use
\begin{equation}
U_{1}(t)=k_{p}(H_{c}-H(t,L))+k_{I}Z-f(t),
\label{defu12}
\end{equation}
where $f(t):=H_{1}(t,L)V_{1}(t,L)/v_{G}$. The boundary conditions \eqref{boundary02} become then
\begin{equation}
\begin{split}
H(t,0)V(t,0)&=Q_{0}(t),\\
H(t,L)V(t,L)&=H_{1}V_{1}(t,L)+v_{G}(1+k_{p})(H(t,L)-H_{c})-v_{G}k_{I}Z,
\label{boundary1}
\end{split}
\end{equation}
where we have actually changed $Z$ and re-define $Z:=Z-k_{p}/k_{I}$, which still satisfies the equation \eqref{Z}. {This can be seen as a feedforward control.}

This new control \eqref{defu12} assumes that $V_{1}(t,L)$ is known at least up to a constant, as $H_{1}(t,L)=H_{c}$ and additional constants can be incorporated into $Z$.
When no knowledge on the target state is available besides $H_{c}$,
it is impossible to stabilize exponentially the system, and the best one can get is the Input-to-State Stability which is given by Theorem \ref{th2}.
However, in the following we will keep working with \eqref{defu12} and \eqref{boundary1} to show Theorem \ref{th1} and the exponential stability of the system, as~the proof of Theorem \ref{th2} and Corollary \ref{cor1}, which uses only the control \eqref{defu1} and \eqref{boundary2}, are easily deduced from the proof of Theorem \ref{th1}.

We introduce the first-order compatibility conditions associated to the boundary conditions \eqref{boundary1} for an initial condition $(H^{0},V^{0},Z^{0})$:
\begin{equation}\label{compat}
\begin{cases}
H^{0}(0)V^{0}(0)=Q_{0}(0),\\
H^{0}(L)V^{0}(L)=H_{1}V_{1}(0,L)+v_{G}(1+k_{p})(H^{0}(L)-H_{c})-k_{I}Z^{0},\\[3pt]
-\partial_{x}\Bigl(H^{0}(0)V^{0}(0)+g\dfrac{gH^{0}(0)^{2}}{2}\Bigr)-(k(V^{0})^{2}(0)-CH^{0}(0))=Q'_{0}(0),\\[8pt]
-\partial_{x}\Bigl(H^{0}(L)V^{0}(L)+g\dfrac{gH^{0}(L)^{2}}{2}\Bigr)-(k(V^{0})^{2}(L)-CH^{0}(L))\\
\hspace*{.5cm}=\partial_{t}(H_{1}V_{1})(0,L)
-v_{G}(1+k_{p})\partial_{x}(H^{0}(L)V^{0}(L))+k_{I}(H^{0}(L)-H_{c}).
\end{cases}
\end{equation}
With such compatibility conditions the system \eqref{Stvenant}, \eqref{Z}, \eqref{boundary1} is well-posed and we have the following theorem due to Wang \cite[Th.\,2.1]{Wang}:
\begin{thm}[Well-posedness]\label{th0}
Let $T>0$, and assume that
$(H_{1},V_{1})$ is well-defined and belongs to $C^{0}([0,T],H^{3}(0,L))$.
There exists $\nu(T)>0$
such that for any $(H^{0},V^{0},Z^{0})\in (H^{2}((0,L))))^{2}\times \mathbb{R}$ satisfying
\begin{equation*}
\lVert H^{0}(\cdot)-H_{1}(0,\cdot)\rVert_{H^{2}(0,L)}+\lVert V^{0}(\cdot)-V_{1}(0,\cdot)\rVert_{H^{2}(0,L)}+|Z^{0}|\leq \nu(T),
\end{equation*}
and satisfying the compatibility conditions \eqref{compat}, the system \eqref{Stvenant}, \eqref{Z}, \eqref{boundary1} has a unique solution $(H,V,Z)\in (C^{0}([0,T],H^{2}((0,L))))^{2}\times C^{1}([0,T])$. Moreover
there exists a positive constant $C(T)$ such that
\begin{multline}
\label{estimate}
\lVert H(t,\cdot)-H_{1}(t,\cdot)\rVert_{H^{2}(0,L)}+\lVert V(t,\cdot)-V_{1}(t,\cdot)\rVert_{H^{2}(0,L)}+|Z|\\
\leq C(T)\bigl(\lVert H^{0}(\cdot)-H_{1}(0,\cdot)\rVert_{H^{2}(0,L)}+\lVert V^{0}(\cdot)-V_{1}(0,\cdot)\rVert_{H^{2}(0,L)}+|Z^{0}|\bigr).
\end{multline}
\end{thm}
To apply the result from \cite{Wang}, note that $Z$ can be seen as a third component of the hyperbolic system with a null propagation speed, a constant initial condition $Z^{0}$, and
$Z(t)$ being thus its value everywhere {on $[0,L]$} including at the boundaries.
\begin{rmk}
\label{r0}
If, in addition, $(H^{0},V^{0})\in H^{3}((0,L);\mathbb{R}^{2})$, then the unique solution $(H,V,Z)$ given by Theorem \ref{th0} belongs to $C^{0}([0,T],H^{3}((0,L);\mathbb{R}^{2}))\times C^{2}([0,T])$ and there exists
a constant $C(T)$ such that
\begin{multline}
\lVert H(t,\cdot)-H_{1}(t,\cdot)\rVert_{H^{3}{(0,L)}}+\lVert V(t,\cdot)-V_{1}(t,\cdot)\rVert_{H^{3}{(0,L)}}+|Z|\\
\leq C(T)\bigl(\lVert H^{0}(\cdot)-H_{1}(0,\cdot)\rVert_{H^{3}{(0,L)}}+\lVert V^{0}(\cdot)-V_{1}(0,\cdot)\rVert_{H^{3}{(0,L)}}+|Z^{0}|\bigr).
%\label{estimate2}
\end{multline}
\end{rmk}
We recall the definition of {(local)} exponential stability.
\begin{defn}
\label{defexpstab}
We say that a trajectory $(H_{1},V_{1})$ is {locally} exponentially stable for the $H^{2}$ norm if there exists $\nu>0$, $C>0$ and $\gamma>0$ such that for any $T>t_{0}\geq 0$ and any $(H^{0},V^{0},Z^{0})$ satisfying
\begin{equation*}
\lVert H^{0}(\cdot)-H_{1}(t_{0},\cdot)\rVert_{H^{2}(0,L)}+\lVert V^{0}(\cdot)-V_{1}(t_{0},\cdot)\rVert_{H^{2}(0,L)}+|Z^{0}|\leq \nu,
\end{equation*}
and the compatibility conditions \eqref{compat}, the system \eqref{Stvenant}, \eqref{Z}, \eqref{boundary1} {with initial condition $(H^{0},V^{0},Z^{0})$ at $t_{0}$} has a unique solution
\[
(H,V,Z)\in \bigl(C^{0}([t_{0},T],H^{2}((0,L)))\bigr)^{2}\times C^{1}([t_{0},T])
\]
and
\begin{multline*}
\lVert H(t,\cdot)-H_{1}(t,\cdot)\rVert_{H^{2}(0,L)}+\lVert V(t,\cdot)-V_{1}(t,\cdot)\rVert_{H^{2}(0,L)}+|Z|\\
\leq Ce^{-\gamma t}\bigl(\lVert H^{0}(\cdot)-H_{1}(t_{0},\cdot)\rVert_{H^{2}(0,L)}+\lVert V^{0}(\cdot)-V_{1}(t_{0},\cdot)\rVert_{H^{2}(0,L)}+|Z^{0}|\bigr),\\ \forall t\in[t_{0},T].
\end{multline*}
\end{defn}
\begin{rmk}
From \eqref{target} and Sobolev inequality, this exponential stability implies in particular the {(local)} exponential convergence of $H(t,L)$ to $H_{c}$.
\end{rmk}

We can now state the main results of this article.

\begin{thm}[Exponential stability]\label{th1}
There exists $\delta>0$ such that, if
\[
\lVert\partial_{t}Q_{0}\rVert_{C^{3}([0,+\infty))}\leq \delta,
\]
then the trajectory $(H_{1},V_{1})$ given by \eqref{target} of system \eqref{Stvenant}, \eqref{Z}, \eqref{boundary1} is exponentially stable for the $H^{2}$ norm if:
\begin{equation}
\begin{split}
&k_{p}>-1\text{ and }k_{I}>0,\\
\text{or}\quad&k_{p}<-1-\frac{gH_{1}(t,L)-V^{2}_{1}(t,L)}{v_{G}V_{1}(t,L)}\;\;\text{and}\;\;k_{I}<0.
\end{split}
\label{cond}
\end{equation}
\end{thm}
This result is proved in Section \ref{s2}.
The main idea of the proof consist in finding a local convex and dissipative entropy for the system \eqref{Stvenant}, \eqref{Z}, \eqref{boundary1}.

In particular, in the case where $Q_{0}$ is constant, we can use the static boundary control \eqref{defu1}, and we have the following corollary:
\begin{cor}\label{cor1}
If $Q_{0}$ is constant, then the steady-state $(H^{*},V^{*})$ of the system \eqref{Stvenant}, \eqref{Z}, \eqref{boundary2} given by \eqref{targetsteady0}--\eqref{targetsteady}
is exponentially stable for the $H^{2}$ norm if:
\begin{equation}
\label{cond1}
\begin{split}
&k_{p}>-1\text{ and }k_{I}>0,\\
\text{or}\quad&k_{p}<-1-\frac{gH^{*}(L)-V^{*2}(L)}{v_{G}V^{*}(L)}\;\;\text{and}\;\;k_{I}<0.
\end{split}
\end{equation}
\end{cor}
\begin{proof}
This is a particular case of Theorem \ref{th1}. To see this, note, as mentioned earlier, that when $Q_{0}$ is constant, then $(H_{1},V_{1})=(H^{*},V^{*})$.
Then, observe that $f(t)$ given in \eqref{defu12} is a constant that can be added in $Z$ (\ie we can re-define $Z:=Z-f(t)$, which still satisfies \eqref{Z}).
\end{proof}

\subsection{Comparison with existing results and contribution of this paper}
Many results exist in the literature concerning this stabilization problem (\eg \cite{dos2008, XuSallet, BastinCoron2013, XuSallet2014, Trinh2015, BastinCoronTamasoiu2015, BastinCoronPI}). To our knowledge the most advanced result for the full non-linear system is \cite{BastinCoronPI} where the authors show that if there is no slope, \ie $C(x) = 0$, then the system can always be stabilized by the PI control \eqref{defu1} as long as the steady-state exists, and they give the sufficient condition
\begin{equation*}%\label{condbastincoron}
k_{p}>0\;\text{ and }\;k_{I}>0.
\end{equation*}
Note that this is the first result that allows an arbitrary size of source term and length.
In this paper, using a different type of Lyapunov function, we manage to show a more general result. Our main contributions are the following:
\begin{itemize}
\item The result holds for an arbitrary friction and also an arbitrary slope $C(x)\in C^{2}([0,L])$. Physically this means that the source can be non-dissipative and increase the energy of the system compared to the case where there is only friction.
\item We find a less restrictive stability condition
\begin{equation*}
k_{p}>-1\;\text{ and }\;k_{I}>0,
\end{equation*}
and we also show that another condition is sufficient:
\begin{equation*}
k_{p}<-1-\frac{gH^{*}(L)-V^{*2}(L)}{v_{G}V^{*}(L)}\;\text{ and }\; k_{I}<0.
\end{equation*}
This one is counter intuitive as $k_{p}<-1$ and $k_{I}<0$. {{It} means that if the height of the water is too high at $L$ the control would reduce the aperture of the gate in $L$ and reduce the flow that we let exit the system,
which intuitively should increase even more the height of the water at $L$.}
\item Our result holds also when stabilizing a slowly varying trajectory rather than a steady-state (that might not exist in practical case). In this case we use a kind of feedforward term in the boundary control {(see \eqref{defu12})}.
\item In addition to the exponential stability, we show the Input-to-State Stability with respect to an unknown inflow. In this case the only knowledge required on the system is the height of the water at $x=L$.
\end{itemize}

{Note that, just like \cite{BastinCoronPI}, this approach uses very little knowledge of the state of the system, as we only measure the height at the boundary $x=L$.}

\begin{rmk}
{When the system is homogeneous, our conditions \eqref{cond1} are optimal (necessary and sufficient) \cite{BastinCoronTamasoiu2015}, \cite[\S 2.2.4.1]{BastinCoron1D}.}
\end{rmk}

\begin{rmk}[Alternative notation in literature]
In the literature, results about PI control of the Saint-Venant equations sometimes leave the step of modeling the spillway, and use a generic formulation of the PI control on the outflow rate of the form\vspace*{-3pt}
\begin{equation*}
H(t,L)V(t,L)=k_{{p}}(H(t,L)-H_{c})-k_{{I}}Z,
\end{equation*}
where $Z$ is the integral term, still given by \eqref{Z}.
Note that, with these notations, the sufficient condition of Corollary \ref{cor1} becomes\vspace*{-3pt}
\begin{equation*}
\begin{split}
&k_{p}>0\;\;\text{and}\;\;k_{I}>0,\\
\text{or}\quad&k_{p}<-\frac{gH^{*}(L)-V^{*2}(L)}{V^{*}(L)}\;\;\text{and}\;\;k_{I}<0.
\end{split}
\end{equation*}
\end{rmk}

\subsection{Case of time-varying input disturbance $Q_{0}(t)$: ISS estimate}
In practical situations, however, we may also have only
little knowledge of the target trajectory $(H_{1},V_{1})$ or the input disturbance $Q_{0}(t)$ and we only know $H_{c}$. In this case we cannot use a controller of the form \eqref{boundary1},
but only a static controller of the form \eqref{boundary2}, namely\vspace*{-3pt}
\begin{equation*}
H(t,L)V(t,L)=v_{G}(1+k_{p})H(t,L)-v_{G}k_{p}H_{c}-v_{G}k_{I}Z.
%\label{boundary22}
\end{equation*}
In this case, it is impossible to aim at stabilizing the target trajectory $(H_{1},V_{1})$, but we still have the Input-to-State Stability with respect to the input disturbance $\partial_{t}Q_{0}$,
\begin{thm}
\label{th2}
There exists $\nu>0$, $\delta>0$, $\gamma>0$ and $C$, such that if\vspace*{-3pt}
\[
\lVert\partial_{t}Q_{0}\rVert_{C^{2}([0,+\infty))}\leq\delta,
\]
then for any $T>0$ and $(H^{0},V^{0})\in {(H^{2}(0,L))^{2}}$ such that\vspace*{-3pt}
\begin{equation*}
\lVert H^{0}-H^{*}\rVert_{H^{2}(0,L)}+\lVert V^{0}-V^{*}\rVert_{H^{2}(0,L)}\leq \nu,
\end{equation*}
the system \eqref{Stvenant}, \eqref{Z}, \eqref{boundary2} with initial condition $(H^{0},V^{0})$ has a unique solution $(H,V)\in C^{0}([0,T],H^{2}(0,L))$ which satisfies the following ISS inequality
\begin{multline}
\label{ISSHV}
\lVert H(t,\cdot)-H_{0}(t,\cdot)\rVert_{H^{2}(0,L)}+\lVert V(t,\cdot)-V_{0}(t,\cdot)\rVert_{H^{2}(0,L)}\\
\leq Ce^{-\gamma t}\biggl(\lVert H^{0}-H^{*},V^{0}-V^{*}\rVert_{H^{2}(0,L)}\\[-3pt]
+\int_{0}^{t}(|\partial_{t}Q_{0}(s)|+|\partial_{tt}^{2}Q_{0}(s)|+|\partial_{ttt}^{3}Q_{0}(s)|)e^{\gamma s}ds\biggr).
\end{multline}
\end{thm}
The proof is given in Appendix \ref{ISS2} and is a consequence from the proof of Theorem~\ref{th1}.
In Section \ref{s2} we {give a few tools to prepare the proof of} Theorem \ref{th1}.
\section{Exponential stability for the \texorpdfstring{$H^{2}$}{H2} norm}
\label{s2}
This section is divided in {two} parts. First we transform the system through a change of variables.
Then we state two lemma, which simplify the proof of Theorem~\ref{th1}.
{We~will then prove Theorem \ref{th1} in Section \ref{sec:proofth1}.}

\subsection{A change of variables}
For any solution of \eqref{Stvenant}, \eqref{Z}, \eqref{boundary1}
we define the perturbations as
\begin{equation}
\begin{pmatrix}
h\\v
\end{pmatrix}=
\begin{pmatrix}
H-H_{1}\\
V-V_{1}
\end{pmatrix}.
\label{change1}
\end{equation}
Let us assume that there exists $\nu\in(0,\nu_{0})$ to be selected later on, such that
\begin{equation*}
\lVert H^{0}(\cdot)-H_{1}(0,\cdot)\rVert_{H^{2}(0,L)}+\lVert V^{0}(\cdot)-V_{1}(0,\cdot)\rVert_{H^{2}(0,L)}+|Z^{0}|\leq \nu.
\end{equation*}
The boundary conditions \eqref{boundary1} can be written in the following form
\begin{equation}
\begin{split}
v(t,0)&=\mathcal{B}_{1}(h(t,0),t),\\
v(t,L)&=\mathcal{B}_{2}(h(t,L),Z,t),
\end{split}
\label{bound1}
\end{equation}
with
\begin{equation}
\begin{split}
\partial_{1}\mathcal{B}_{1}(0,t)&=-\frac{V_{1}(t,0)}{H_{1}(t,0)},\\
\partial_{1}\mathcal{B}_{2}(0,0,t)&=\frac{v_{G}(1+k_{p})-V_{1}(t,L)}{H_{1}(t,L)},\\
\partial_{2}\mathcal{B}_{2}(0,0,t)&=-\frac{v_{G}k_{I}}{H_{1}(t,L)}.
\end{split}
\label{k1k30}
\end{equation}
We introduce the following change of variables:
\begin{equation}
\boldsymbol{u}: =\begin{pmatrix}
u_{1}\\u_{2}
\end{pmatrix}=
\begin{pmatrix}
v+\sqrt{\sfrac{g}{H_{1}}}\,h\\[2pt]
v-\sqrt{\sfrac{g}{H_{1}}}\,h
\end{pmatrix}.
\label{change2}
\end{equation}
Note that this change of variables is very similar to the change of variables used in \hbox{\cite{BastinCoron22,HS}} with the only difference that $(H_{1},V_{1})$ is not a steady-state anymore.
It corresponds to the transformation in Riemann coordinates for the perturbations.
Indeed, denoting $S$, $F$ and $G$ by
\begin{gather}
S(x,t)=\begin{pmatrix}\sqrt{\sfrac{g}{H_{1}(t,x)}} & 1\\[3pt]
-\sqrt{\sfrac{g}{H_{1}(t,x)}} & 1\end{pmatrix},
\label{defS}\\
F\begin{pmatrix}H\\V\end{pmatrix}=\begin{pmatrix}
V & H\\
g & V
\end{pmatrix},
\label{defF}\qquad
G\begin{pmatrix}H\\V\end{pmatrix}=\begin{pmatrix}0\\\sfrac{kV^{2}}{H}-C(x)\end{pmatrix},
\end{gather}
and using \eqref{Stvenant}, \eqref{Z}, \eqref{boundary1}, \eqref{target}, \eqref{change1}--\eqref{change2}, one has
\begin{equation}
\label{sys1}
\begin{split}
\partial_{t}u_{1}+\lambda_{1}(\boldsymbol{u},x,t)\partial_{x}u_{1}+l_{1}(\boldsymbol{u},x,t)\partial_{x}u_{2}+B_{1}(\boldsymbol{u},x,t)&=0,\\
\partial_{t}{u_{2}}-\lambda_{2}(\boldsymbol{u},x,t)\partial_{x}u_{2}+l_{2}(\boldsymbol{u},x,t)\partial_{x}u_{1}+B_{2}(\boldsymbol{u},x,t)&=0,
\end{split}
\end{equation}
where\vspace*{-5pt}
\begin{multline}
\label{defA}
\begin{pmatrix}\lambda_{1}(\boldsymbol{u},x,t) & l_{1}(\boldsymbol{u},x,t)\\l_{2}(\boldsymbol{u},x,t) & -\lambda_{2}(\boldsymbol{u},x,t)\end{pmatrix}\\[-5pt]
=S(x,t)F\left(S^{-1}(x,t)\boldsymbol{u}+\begin{pmatrix}H_{1}(t,x)\\V_{1}(t,x)\end{pmatrix}\right)S^{-1}(x,t)=:A(\boldsymbol{u},x,t),
\end{multline}
\par\noindent
\vspace*{-5pt}and\vspace*{-5pt}\enlargethispage{\baselineskip}%
\begin{multline}
\label{defB}
B(\boldsymbol{u},x,t)=\begin{pmatrix}B_{1}(\boldsymbol{u},x,t)\\B_{2}(\boldsymbol{u},x,t)\end{pmatrix}
\\
\shoveright{=S(x,t)F\biggl(S^{-1}(x,t)\boldsymbol{u}+\begin{pmatrix}H_{1}(t,x)\\V_{1}(t,x)\end{pmatrix}\biggr)\biggl(\begin{pmatrix} \partial_{x}H_{1}(t,x)\\ \partial_{x}V_{1}(t,x)\end{pmatrix}+\partial_{x}(S^{-1})\boldsymbol{u}\biggr)}\\
\shoveright{+S\partial_{t}\begin{pmatrix}H_{1}(t,x)\\V_{1}(t,x)\end{pmatrix}+S(x,t)G\biggl(S^{-1}(x,t)\boldsymbol{u}+\begin{pmatrix}H_{1}(t,x)\\V_{1}(t,x)\end{pmatrix}\biggr)}\\
-\partial_{t}S(x,t)S^{-1}(x,t)\boldsymbol{u}.
\end{multline}
\par\noindent
\vspace*{-3pt}Therefore,\vspace*{-2pt}
\begin{gather}
\begin{aligned}
\lambda_{1}(0,x,t)&=V_{1}+\sqrt{gH_{1}},&\lambda_{2}(0,x,t)&=\sqrt{gH_{1}}-V_{1}\label{Lambda},\\
l_{1}(0,x,t)&=B_{1}(\boldsymbol{0},x,t)=0,& l_{2}(0,x,t)&=B_{2}(\boldsymbol{0},x,t)=0,
\end{aligned}
\\
\begin{split}
\frac{\partial{B}_{1}}{\partial_{u}}(0,x,t)&=\gamma_{1}(t,x)u_{1}(t,x)+\gamma_{2}(t,x)u_{2}(t,x),\\
\frac{\partial{B}_{2}}{\partial_{u}}(0,x,t)&=\delta_{1}(t,x)u_{1}(t,x)+\delta_{2}(t,x)u_{2}(t,x),
\end{split}
\label{l1l2}
\end{gather}
where\vspace*{-7pt}
\begin{equation}
\begin{split}
\gamma_{1}&=\frac{3}{4}\sqrt{\sfrac{g}{H_{1}}}\,H_{1x}+\frac{3}{4}V_{1x}+\frac{kV_{1}}{H_{1}}-\frac{kV_{1}^{2}}{2H_{1}^{2}}\sqrt{\sfrac{H_{1}}{g}},\\
\gamma_{2}&=\frac{1}{4}\sqrt{\sfrac{g}{H_{1}}}\,H_{1x}+\frac{1}{4}V_{1x}+\frac{kV_{1}}{H_{1}}+\frac{kV_{1}^{2}}{2H_{1}^{2}}\sqrt{\sfrac{H_{1}}{g}},\\
\delta_{1}&=-\frac{1}{4}\sqrt{\sfrac{g}{H_{1}}}\,H_{1x}+\frac{1}{4}V_{1x}+\frac{kV_{1}}{H_{1}}-\frac{kV_{1}^{2}}{2H_{1}^{2}}\sqrt{\sfrac{H_{1}}{g}},\\
\delta_{2}&=-\frac{3}{4}\sqrt{\sfrac{g}{H_{1}}}\,H_{1x}+\frac{3}{4}V_{1x}+\frac{kV_{1}}{H_{1}}+\frac{kV_{1}^{2}}{2H_{1}^{2}}\sqrt{\sfrac{H_{1}}{g}}.
\end{split}
\label{gamma}
\end{equation}
{For} the boundary conditions, there exists $\nu_{1}\in(0,\nu_{0})$ such that for any $\nu\in(0,\nu_{1})$, one has:\vspace*{-5pt}
\begin{equation}
\begin{split}
u_{1}(t,0)&=\mathcal{D}_{1}(u_{2}(t,0),t),\\
u_{2}(t,L)&=\mathcal{D}_{2}(u_{1}(t,L),Z,t),\\
\dot Z&=\frac{(u_{1}(t,L)-u_{2}(t,L))}{2}\sqrt{\sfrac{H_{1}(t,L)}{g}},
\end{split}
\label{bound}
\end{equation}
where $\mathcal{D}_{1}$ and $\mathcal{D}_{2}$ are $C^{2}$ functions and\vspace*{-5pt}
\begin{equation}
\begin{split}
\partial_{1}\mathcal{D}_{1}(0,t)&=-\frac{\lambda_{2}(0)}{\lambda_{1}(0)},\\
\partial_{1}\mathcal{D}_{2}(0,0,t)&=-\frac{\lambda_{1}(L)-v_{G}(1+k_{p})}{\lambda_{2}(L)+v_{G}(1+k_{p})},\\
\partial_{2}\mathcal{D}_{2}(0,0,t)&=-2\frac{v_{G}k_{I}\sqrt{\sfrac{g}{H_{1}(t,L)}}}{v_{G}(1+k_{p})+\lambda_{2}(t,L)}.
\end{split}
\label{k1k3}
\end{equation}
Expression \eqref{gamma} is simply a computation, very similar to what is done in \cite{HS} for instance, while the derivation of \eqref{bound} and \eqref{k1k3} are detailed in the appendix. In~the following we denote for simplicity
\[
k_{2}:=\partial_{1} \mathcal{D}_{1}(0,t),\quad k_{1}:=\partial_{1}\mathcal{D}_{2}(0,0,t)\quad\text{and}\quad k_{3}:=-\partial_{2}\mathcal{D}_{2}(0,0,t),
\]
\begin{rmk}
Obviously, from the change of variables \eqref{change1}--\eqref{change2}, the exponential stability of
the system \eqref{Stvenant}, \eqref{Z}, \eqref{boundary1}
is equivalent to the exponential stability of the steady-state $\boldsymbol{u}^{*}=0$ for the system \eqref{sys1}, \eqref{bound}.
\label{r1}
\end{rmk}
As the operator $A$, given by \eqref{defA}, is a {$C^{2}$ function in $\boldsymbol{u}$,
$t$ and $x$ (and {in particular~$C^{1}$})}
and as, from \eqref{l1l2} and \eqref{fluvial}, $\lambda_{1}(\boldsymbol{0},x,t)>0>-\lambda_{2}(\boldsymbol{0},x,t)$,
there exists $\nu_{2}\in(0,\nu_{1})$ {depending only on $H_{\max}$, $\alpha$}
and $E\in C^{1}(\mathcal{B}_{\nu_{2}}\times(0,L)\times[0,+\infty);\mathcal{M}_{2}(\mathbb{R}))$, where
$\mathcal{B}_{\nu_{2}}\subset\mathbb{R}^{2}$ is the disc of radius $\nu_{2}$ and center $0$,
such that
for any $\lVert\boldsymbol{u}(t,\cdot)\rVert_{H^{2}(0,L)}\leq \nu_{2}$,
\begin{equation}
\begin{split}
E(\boldsymbol{u}(t,x),x,t)A(\boldsymbol{u}(t,x),x,t)&=D(\boldsymbol{u}(t,x),x,t)E(\boldsymbol{u}(t,x),x,t),\\
E(\boldsymbol{0},x,t)&=\mathrm{Id},
\end{split}
\label{definv}
\end{equation}
where $D(\boldsymbol{u}(t,x),x,t)=(D_{i}(\boldsymbol{u}(t,x),x,t))_{i\in{1,2}}$ is a diagonal matrix and $\mathrm{Id}$ is the identity matrix.
Before going any further, let us
note a few useful properties of these functions. {For simplicity in the following we will denote for any $n\in\mathbb{N}^{*}$ and any function $U\in L^{\infty}((0,T)\times(0,L);\mathbb{R}^{n})$ (resp. $L^{\infty}((0,L);\mathbb{R}^{n})$)
\begin{equation*}
\begin{split}
\lVert U\rVert_{\infty}&:=\lVert U\rVert_{L^{\infty}((0,T)\times(0,L);\mathbb{R}^{n})},\\
(\resp\lVert U\rVert_{\infty}&:=\lVert U\rVert_{L^{\infty}((0,L);\mathbb{R}^{n})}).
\end{split}
\end{equation*}
We may also denote $\lVert\boldsymbol{u}\rVert_{H^{2}(0,L)}$ instead of $\lVert\boldsymbol{u}(t,\cdot)\rVert_{H^{2}(0,L)}$ to lighten the expressions.}
From the definition of $A$ given in \eqref{defA},
and from \eqref{fluvial},
for $\lVert\boldsymbol{u}\rVert_{H^{2}(0,L)}\leq \nu_{2}$, there exists a constant $C_{1}$ depending only on $H_{\max}$, $\alpha$ and $\nu_{2}$ such that we have the following estimates
\begin{equation}
\label{ineq}
\begin{split}
\max\bigl(\lVert \partial_{t}(A(\boldsymbol{u}(t,x),x,t)-A(\boldsymbol{0},x,t))\rVert_{\infty},\lVert \partial_{t}(D(\boldsymbol{u}(t,x),x,t)-D(\boldsymbol{0},x,t))\rVert_{\infty},&\\
\lVert \partial_{t}(E(\boldsymbol{u}(t,x),x,t))\rVert_{\infty}\bigr)&\\
\leq C_{1}\bigl(\lVert\boldsymbol{u}\rVert_{\infty}(\lVert\partial_{t}H_{1}\rVert_{\infty}+\lVert\partial_{t}V_{1}\rVert_{\infty})+\lVert\partial_{t}\boldsymbol{u}\rVert_{\infty}\bigr),&\\
\max\bigl(\lVert \partial_{{x}}(A(\boldsymbol{u}(t,x),x,t)-A(\boldsymbol{0},x,t))\rVert_{\infty},\lVert \partial_{{x}}(D(\boldsymbol{u}(t,x),x,t)-D(\boldsymbol{0},x,t)),&\\
\lVert \partial_{{x}}(E(\boldsymbol{u}(t,x),x,t))\rVert_{\infty}\bigr)&\\
\leq C_{1} \bigl(\lVert\boldsymbol{u}\rVert_{\infty}(\lVert\partial_{x}H_{1}\rVert_{\infty}+\lVert\partial_{x}V_{1}\rVert_{\infty})+\lVert\partial_{x}\boldsymbol{u}\rVert_{\infty} \bigr).&
\end{split}
\end{equation}
For $E$ and $D$, this comes from the fact that $E$ and $D$ are $C^{\infty}$ functions with respect to the coefficients of $A$ ({recall} that $D$ is the matrix of eigenvalues of $A$), and that $A\in C^{2}(\mathcal{B}_{\eta_{0}}; C^{1}([0,+\infty)\times[0,L]))$.

\subsection{Two useful lemmas}
We introduce now two lemma, which simplify the proof of Theorem \ref{th1}. The first one is a classical result about Lyapunov functions.
\begin{lem}
Let $V: (H^{2}(0,L))^{2}\times\mathbb{R}\times\mathbb{R}_{+}\to \mathbb{R}_{+}^{*}$, {$\eta>0$,} and $c>0$ such that for any $(U,z,t)\in (H^{2}(0,L))^{2}\times\mathbb{R}\times\mathbb{R}_{+}$, {with $\|\boldsymbol{U}\|_{H^{2}}+|Z|\leq \eta$,}
\begin{equation}
c\bigl(\lVert \boldsymbol{U} \rVert_{H^{2}(0,L)}^{{2}}+|z|^{{2}}\bigr)\leq V(\boldsymbol{U},z,t)\leq \frac{1}{c} \bigl(\lVert \boldsymbol{U} \rVert_{H^{2}(0,L)}^{{2}}+|z|^{{2}}\bigr).
\label{cV1}
\end{equation}
If {for any $T\geq t_{0}>0$} there exists $\gamma>0$ {independent of $t_{0}$ and $T$,}
and ${\nu}>0$ {depending only on $T-t_{0}$} such that, for
any solution
$(\boldsymbol{u},Z)$ of the system \eqref{sys1}, \eqref{bound} with initial conditions
satisfying $\lVert \boldsymbol{u}({t_{0}},\cdot)\rVert_{H^{2}(0,L)}+|Z({t_{0}})|\leq {\nu}$, {the differential inequality}
\begin{equation}
\frac{d}{dt}\left[V(\boldsymbol{u}(t,\cdot),t)\right]<-\gamma V(\boldsymbol{u}(t,\cdot),t),\quad\forall t\in(t_{0},T),
\label{cV2}
\end{equation}
{holds} in a distribution sense, then the system \eqref{sys1}, \eqref{bound} is exponentially stable for the $H^{2}$ norm and $V$ is called a Lyapunov function for the system \eqref{sys1}, \eqref{bound}.
\label{lem1}
\end{lem}
This first lemma reduces the problem of proving the exponential stability to finding a Lyapunov function $V$ for the system \eqref{sys1}, \eqref{bound}. A proper definition of a differential inequality in a distribution sense as in \eqref{cV2} can be found in \cite{C1}. To lighten this article we do not give a proof of this classical lemma, although a proof for a very similar case (Lyapunov function that does not depend explicitly on time and for the~$C^{1}$ norm instead) can be found for instance in \cite[Prop.\,2.1]{C1}, and is easily extended to this~case. 

\begin{rmk}
Note that $\nu$ may depend on $t_{0}$ and $T$ with this definition, contrarily to the Definition \ref{defexpstab} of exponential stability. However, this is not an issue since we can deduce, thanks to \eqref{cV2}, the existence of $\nu$ independent of $t_{0}$ and $T$ such that \eqref{cV2} holds with $\gamma/2$ instead of $\gamma$. Indeed, assume that the assumption of Lemma \ref{lem1} holds and select $t_{0}=0$, and $T_{1}>0$ such that $e^{-\gamma T_{1}/2}<c^{2}/2$ where $c$ is the positive constant involved in \eqref{cV1} (note that $\gamma$ does not depend on $t_{0}$ and $T_{1}$). From \eqref{cV2} there exists $\nu$ depending on $T_{1}$ (that we denote $\nu(T_{1})$ in the following) such that if $\lVert \boldsymbol{u}({0},\cdot)\rVert_{H^{2}(0,L)}+|Z({0})|\leq {\nu}(T_{1})$, then
\begin{equation*}
V(\boldsymbol{u}(T_{1},\cdot))\leq V(\boldsymbol{u}(0,\cdot))e^{-\gamma T_{1}}.
\end{equation*}
Thus, using \eqref{cV1}
\begin{equation}
\label{eq:esimrmk}
\lVert \boldsymbol{u}(T_{1},\cdot)\rVert_{H^{2}(0,L)}+|Z(T_{1})|\leq c^{2}e^{-\gamma T_{1}}\left(\lVert \boldsymbol{u}({0},\cdot)\rVert_{H^{2}(0,L)}+|Z({0})|\right)\leq \frac{1}{2}\,\nu (T_{1}).
\end{equation}
Since $\nu$ only depends on $T-t_{0}$, \eqref{eq:esimrmk} means that the system with initial condition $(\boldsymbol{u}(T_{1},\cdot)$, $Z(T_{1}))$ has a solution on $[T_{1},2T_{1}]$ and, thanks to \eqref{cV1}--\eqref{cV2} and the choice of $T_{1}$,
\begin{equation*}
\begin{split}
\lVert \boldsymbol{u}(t,\cdot)\rVert_{H^{2}(0,L)}+|Z(t)|&\leq c^{2}e^{-\gamma(t-T_{1})}c^{2}e^{-\gamma T_{1}}\left(\lVert \boldsymbol{u}({0},\cdot)\rVert_{H^{2}(0,L)}+|Z({0})|\right)\\
&\leq c^{2}e^{-\gamma(t-T_{1})}e^{-\gamma T_{1}/2}\left(\lVert \boldsymbol{u}({0},\cdot)\rVert_{H^{2}(0,L)}+|Z({0})|\right)\\
&\leq c^{2}e^{-\gamma t/2}\left(\lVert \boldsymbol{u}({0},\cdot)\rVert_{H^{2}(0,L)}+|Z({0})|\right),\quad\forall t\in[T_{1},2T_{1}]
\end{split}
\end{equation*}
and
\begin{equation*}
\lVert \boldsymbol{u}(2T_{1},\cdot)\rVert_{H^{2}(0,L)}+|Z(2T_{1})|\leq c^{4}e^{-2\gamma T_{1}}\left(\lVert \boldsymbol{u}({0},\cdot)\rVert_{H^{2}(0,L)}+|Z({0})|\right)\leq \frac{1}{4}\,\nu (T_{1}).
\end{equation*}
Hence we can iterate on $[2T_{1}, 3T_{1}],\dots,[nT_{1},(n+1)T_{1}]$, and so on, for any $n\in\mathbb{N}$ and we obtain the exponential stability estimate on $[0,+\infty)$ and in particular on any $[t_{0},T]$ with $\nu := \nu(T_{1})$ that does not depend on $T$ or $t_{0}$.
\end{rmk}

Our second lemma seems very natural:
\begin{lem}\label{lem3}
There exists $l>0$ and $C>0$
such that if $\lVert\partial_{t}Q_{0}\rVert_{C^{3}([0,+\infty))}\leq l$, then
\begin{equation*}
\max\left(\lVert\partial_{t}H_{1}\rVert_{C^{1}([0,+\infty), {C^{0}([0,L])})},\lVert\partial_{t}V_{1}\rVert_{C^{1}([0,+\infty), {C^{0}([0,L])})}\right)< C \lVert\partial_{t}Q_{0}\rVert_{C^{3}([0,+\infty))}.
\end{equation*}
\end{lem}
This is a consequence of the ISS property (Proposition \ref{propISS}) and Remark \ref{rmkHpISS} for $p=3$ and is shown in Appendix \ref{app:prooflem3}.
Thanks to this lemma, we now only need to show Theorem \ref{th1} with a bound on $\partial_{t}H_{1}$ and $\partial_{t}V_{1}$ rather than a bound on $\partial_{t}Q_{0}$.

\section{{Proof of Theorem \ref{th1}}}
\label{sec:proofth1}
From Theorem {\ref{th0}}, Remark \ref{r1}, and Lemma \ref{lem1}, one only needs to find a Lyapunov function $V: (H^{2}(0,L))^{2}\times\mathbb{R}\times\mathbb{R}_{+}\to \mathbb{R}_{+}^{*}$ satisfying \eqref{cV1} and \eqref{cV2}.
{We will proceed as follows: first we introduce a Lyapunov function defined up to two positive functions $f_{1}$ and $f_{2}$ and two positive constants $\mu$ and $q$. Second, we show a differential inequality satisfied by this Lyapunov function candidate with respect to $f_{1}$, $f_{2}$, $\mu$ and~$q$ (Proposition \ref{propf1f2}). Then we give sufficient conditions such that this differential inequality simplifies to \eqref{cV2}. Finally, we show how to choose $f_{1}$, $f_{2}$, and then~$q$ and~$\mu$ such that these sufficient conditions are satisfied, together with \eqref{cV1}.}

Let {us} define the following functional on $H^{2}(0,L)\times\mathbb{R}\times\mathbb{R}_{+}$:
\begin{multline}
\label{Va}
V_{a}(\boldsymbol{U},z,t):=\int_{0}^{L}f_{1}(t,x)e^{-\mu x}(E(\boldsymbol{U}(x),x,t)\boldsymbol{U}{(x)})_{1}^{2}\\
+f_{2}(t,x)e^{\mu x}(E(\boldsymbol{U}(x),x,t)\boldsymbol{U}{(x)})_{2}^{2}dx+qz^{2},
\end{multline}
where $f_{1}$, $f_{2}$ are positive and bounded functions which will be defined later on, and~$\mu$ and $q$ are positives constant which will also be defined later on. {Recall that $E$ is given by \eqref{definv}. We introduce the following candidate Lyapunov function defined for $H^{2}$ trajectories of \eqref{sys1}:
\begin{equation}
\label{def-V-tot}
V(\boldsymbol{u}(t,\cdot),Z(t),t) = V_{a}(\boldsymbol{u}(t,\cdot),Z(t),t)+V_{b}(\boldsymbol{u},t)+V_{c}(\boldsymbol{u},t),
\end{equation}
where
\begin{equation}
\begin{split}
V_{b}(\boldsymbol{u},t) &= V_{a}(\partial_{t}\boldsymbol{u}(t,\cdot),\dot Z(t),t),\\
V_{c}(\boldsymbol{u},t) &= V_{a}(\partial_{t}^{2}\boldsymbol{u}(t,\cdot),\ddot Z(t),t).
\end{split}
\label{VbVc2}
\end{equation}
This functional is a priori only defined for trajectories of \eqref{sys1}, however using \eqref{sys1} its definition can be extended to $H^{2}(0,L)\times\mathbb{R}\times\mathbb{R}_{+}$ as well (see Appendix \ref{sec:VbVc} for more details). We have the following proposition
\begin{prop}
\label{propf1f2}
There exists $\delta>0$ such that if
\begin{equation*}
\max\bigl(\lVert\partial_{t}H_{1}\rVert_{C^{1}([t_{0},\infty);{C^{0}([0,L])})},\lVert\partial_{t}V_{1}\rVert_{C^{1}([t_{0},\infty);{C^{0}([0,L])})}\bigr)<\delta,
\end{equation*}
then for any $T>0$ and $t_{0}\in[0,T)$,
there exists $\nu>0$ {depending only on $T-t_{0}$}
such that for any solution $(\boldsymbol{u},Z)$ of \eqref{sys1} belonging to $C^{0}([t_{0},T],H^{3}(0,L))\times C^{2}([t_{0},T])$ with initial condition $\boldsymbol{u}^{0}$ and $Z^{0}$ at time $t_{0}$ satisfying
\begin{equation*}
\left(\lVert \boldsymbol{u}^{0}\rVert_{H^{2}(0,L)}+|Z^{0}|\right)<\nu,
\end{equation*}
one has the following differential inequalities for any $t\in[t_{0},T]$
\begin{multline}
\label{dVtot}
\frac{dV(\boldsymbol{u}(t,\cdot),Z,t)}{dt}\leq-\mu\min\limits_{x\in[0,L]}(\lambda_{1},\lambda_{2})V(\boldsymbol{u}(t,\cdot),Z,t)\\
+\left[f_{1}\lambda_{1}k_{2}^{2}-\lambda_{2}f_{2}\right]\left(u_{2}^{2}(t,0)+(\partial_{t}u_{2}(t,0))^{2}+(\partial_{tt}^{2}u_{2}(t,0))^{2}\right)\\
-I_{1}(u_{1}(t,L),Z)-I_{1}(\partial_{t}u_{1}(t,L),\dot Z)-I_{1}(\partial_{tt}^{2}u_{1}(t,L),\ddot Z)\\
-\int_{0}^{L}I_{2}((E\boldsymbol{u})_{1},(E\boldsymbol{u})_{2})+I_{2}((E\partial_{t}\boldsymbol{u})_{1},(E\partial_{t}\boldsymbol{u})_{2})+I_{2}((E\partial_{tt}^{2}\boldsymbol{u})_{1},(E\partial_{tt}^{2}\boldsymbol{u})_{2})dx\\
+C\bigl(\lVert \boldsymbol{u}\rVert_{\infty}+\lVert \partial_{x} \boldsymbol{u}\rVert_{\infty} \bigr) \Bigl(\lVert\boldsymbol{u} \rVert_{L^{2}(0,L)}^{2}+\lVert\partial_{t}\boldsymbol{u} \rVert_{L^{2}(0,L)}^{2}
+\lVert\partial_{tt}^{2}\boldsymbol{u} \rVert_{L^{2}(0,L)}^{2}\\
\shoveright{+\bigl(\lVert \boldsymbol{u}\rVert_{\infty}+\lVert \partial_{x} \boldsymbol{u}\rVert_{\infty}\bigr)^{{2}}
+|u_{2}(t,0)^{2}|+(|u_{1}(t,L)|+|Z|)^{2}}\\
\shoveright{+|\partial_{t}u_{2}(t,0)^{2}|+(|\partial_{t}u_{1}(t,L)|+|\dot Z|)^{2}
+|\partial_{tt}^{2}u_{2}(t,0)^{2}|
+(|\partial_{tt}^{2}u_{1}(t,L)|+|\ddot Z|)^{2}\Bigr)}\\
+C\delta\left(|u_{2}(t,0)|^{2}+(|u_{1}(t,L)|+|Z|)^{2}+|\partial_{t}u_{2}(t,0)|^{2}+(|\partial_{t}u_{1}(t,L)|+|\dot Z|)^{2}\right)+C\delta V,
\end{multline}
where $C$ is a constant independent of $t_{0}$ and $T$, and $I_{1}$, $I_{2}$ denote the quadratic forms given by
\begin{equation*}
\begin{split}
I_{1}(x,y)&=\left(\lambda_{1}f_{1}(L)e^{-\mu L}-\lambda_{2}f_{2}(L)e^{\mu L}k_{1}^{2}\right)x^{2}\\
&\hspace*{2.06cm}+{\Bigl(q\sqrt{\sfrac{H_{1}}{g}}\,k_{3}-\lambda_{2}f_{2}(L)e^{\mu L}k_{3}^{2}-\mu\min\limits_{x\in[0,L]}(\lambda_{1},\lambda_{2})q \Bigr)}y^{2}\\
&\hspace*{2.06cm}+\bigl(2\lambda_{2}f_{2}(L)e^{\mu L}k_{3}k_{1}-{q\sqrt{\sfrac{H_{1}}{g}}}\,(k_{1}-1)\bigr)xy,\\
I_{2}(x,y)&=\bigl((-\lambda_{1}f_{1})_{x}+2f_{1}\gamma_{1}(t,x)-\partial_{t}f_{1}\bigr)e^{-\mu x}x^{2}\\
&\hspace*{2.06cm}+\bigl((\lambda_{2}f_{2})_{x}+2f_{2}\delta_{2}(t,x)-\partial_{t}f_{2}\bigr)e^{\mu x}y^{2}\\
&\hspace*{2.06cm}+2 \bigl(\gamma_{2}f_{1}e^{-\mu x}+\delta_{1}f_{2}e^{\mu x} \bigr)xy.
\end{split}
%\label{I1I2}
\end{equation*}
\end{prop}
This proposition is showed in Appendix \ref{sec:propf1f2}. We can now use this to derive sufficient conditions for the Lyapunov function candidate to satisfy \eqref{cV1} and \eqref{cV2} and prove Theorem \ref{th1}.}

\Subsection{Sufficient conditions for a Lyapunov function}
\label{sec:sufficient}
Let $T>{t_{0}\geq0}$ and $(\boldsymbol{u}^{0},Z^{0})\in H^{2}(0,L)\times\mathbb{R}$ satisfying the compatibility condition \eqref{compat} and such that
\begin{equation}
\label{boundu1}
\left(\lVert \boldsymbol{u}^{0}\rVert_{H^{2}(0,L)}+|Z^{0}|\right)<\nu,
\end{equation}
where $\nu$ is a constant to be chosen later on but such that $\nu<\min(\nu_{2},\nu(T-t_{0}))$. Recall that ${T\to \nu(T)}$ is given by Theorem \ref{th0} {and $\nu_{2}$ is chosen such that \eqref{definv} holds}.

From Theorem \ref{th0} there exists a unique solution $\boldsymbol{u}\in C^{0}([{t_{0}},T],H^{2}(0,L))$.
To~obtain this, set
\begin{align*}
(\widetilde{H}_{1}(t,\cdot), \widetilde{V}_{1}(t,\cdot)) &= (H_{1}(t-t_{0},\cdot), V_{1}(t-t_{0},\cdot)), \\
(\widetilde{H}(t,\cdot), \widetilde{V}(t,\cdot)) &= (H(t-t_{0},\cdot), V(t-t_{0},\cdot)),\\
\widetilde{Z} &= Z(t-t_{0}). 
\end{align*}
Since $(\widetilde{H}_{1}, \widetilde{V}_{1})$ is still in $C^{0}([0,T-t_{0}], H^{3}(0,L))$ (besides, note that it satisfies the same upper bound in the $H^{2}$ norm as $(H_{1}, V_{1})$), Theorem \ref{th0} still applies and there exists a unique solution $(\widetilde{H},\widetilde{V})$ in $C^{0}([0,T-t_{0}],H^{2}(0,L))$ satisfying \eqref{estimate} on $[0,T-t_{0}]$ with $(\widetilde{H}_{1}, \widetilde{V}_{1})$ instead of $(H_{1}, V_{1})$, provided that
\begin{equation*}
\lVert \widetilde{H}(0,\cdot)-\widetilde{H}_{1}(0,\cdot)\rVert_{H^{2}(0,L)}+\lVert \widetilde{V}(0,\cdot)-\widetilde{V}_{1}(0,\cdot)\rVert_{H^{2}(0,L)}+|Z(t_{0})|\leq \nu(T-t_{0}),
\end{equation*}
which is exactly \eqref{boundu1}. Thus $(H(t,\cdot), V(t,\cdot)) = (\widetilde{H}(t+t_{0},\cdot), \widetilde{V}(t+t_{0},\cdot))$ belongs to $C^{0}([t_{0},T],H^{2}(0,L))$.
In order to use Proposition \ref{propf1f2}, we suppose in addition that $(\boldsymbol{u}^{0},Z^{0})\in H^{3}(0,L){\times\mathbb{R}}$, and that \eqref{boundu1} also holds for the $H^{3}$ norm instead of the~$H^{2}$ norm in $u$. From Remark \ref{r0},
$(\boldsymbol{u},Z)\!\in\! C^{0}([{t_{0}},T]\times H^{3}(0,L))\times C^{3}([{t_{0}},T])$. This assumption will be later relaxed later on by density. From Lemma \ref{lem3}, instead of assuming a bound on $\|\partial_{t} Q_{0}\|_{C^{3}([0,+\infty)}$ we can assume that
\begin{equation}
\max\bigl(\lVert\partial_{t}H_{1}\rVert_{C^{1}([t_{0},\infty);{C^{0}([0,L])})},\lVert\partial_{t}V_{1}\rVert_{C^{1}([t_{0},\infty);{C^{0}([0,L])})}\bigr)<\delta,
\label{boundHV}
\end{equation}
where $\delta$ is a positive constant {independent of $T$}. Let $(f_{1}$, $f_{2})\in C^{1}([0,L];(0,+\infty))$ and $q>0$, $\mu>0$ to be defined later on. From Proposition \ref{propf1f2}, there exists~$\nu_{1}$ and~$\delta_{1}$ such that if $\delta<\delta_{1}$, then the differential inequality \eqref{dVtot} holds.
In the expression of \eqref{dVtot}, one can see that three identical quadratic forms appear in the integral in $((E\partial_{t}^{i}\boldsymbol{u})_{1},(E\partial_{t}^{i}\boldsymbol{u})_{2}),$ $i=0,1,2$, as well as three identical quadratic forms at the boundaries in $(\partial_{t}^{i}u_{1}(t,L),\partial_{t}^{i}Z),$ $i=0,1,2$, and three identical terms proportional respec\-tively to $(\partial_{t}^{i}u_{2}(t,0)),$ $i=0,1,2$.
Thus a sufficient condition {so that there exists $\mu>0$ such} that $V$ is strictly decreasing would be that the square terms and the forms that appear at the boundaries are negative-definite and the quadratic form in the integral is negative-definite, \ie the three following conditions:
\begin{enumerate}
\item \emph{Condition at $0$}
\begin{equation}
\label{c1}
\frac{\lambda_{2}f_{2}(0)}{\lambda_{1}f_{1}(0)}>k_{2}^{2}.
\end{equation}
\item \emph{Condition at $L$}
\begin{subequations}
\label{c2}
\begin{align}
\label{c2a}
&\frac{\lambda_{1}f_{1}(L)}{\lambda_{2}f_{2}(L)}>k_{1}^{2},\\
\label{c2b}
&\left(\lambda_{1}f_{1}(L)-\lambda_{2}f_{2}(L)k_{1}^{2}\right)\left(q\sqrt{\sfrac{H_{1}}{g}}-\lambda_{2}f_{2}(L)k_{3}\right)k_{3}\\
&\hspace*{3cm}-\Bigl(\lambda_{2}f_{2}(L)k_{3}k_{1}-\frac{1}{2}q\sqrt{\sfrac{H_{1}}{g}}\left(k_{1}-1\right)\Bigr)^{2}>0. \notag
\end{align}
\end{subequations}
\item \emph{Condition from the integral}
\begin{equation}\label{c3}
\begin{aligned}
&\bigl((-\lambda_{1}f_{1})_{x}+2f_{1}\gamma_{1}(t,x)-\partial_{t}f_{1}\bigr)>0,
\\
&\bigl((-\lambda_{1}f_{1})_{x}+2f_{1}\gamma_{1}(t,x)-\partial_{t}f_{1}\bigr)\bigl((\lambda_{2}f_{2})_{x}+2f_{2}\delta_{2}(t,x)-\partial_{t}f_{2}\bigr)
\\
& \hspace*{4.2cm}-\bigl(\gamma_{2}f_{1}+\delta_{1}f_{2}\bigr)^{2}>0,\quad\forall (t,x)\in [t_{0},T]\times(0,L).
%\label{c3b}
\end{aligned}
\end{equation}
\end{enumerate}
Let assume for the moment that \eqref{c1}--\eqref{c3} are satisfied
for any $\delta\in(0,\delta_{{2}})$ where $\delta_{{2}}$ is a positive constant. {We are going to show that \eqref{cV2} hold.}
Then, as the inequalities \eqref{c1}--\eqref{c3} are strict, by continuity there exist $\mu>0$ such that the square terms and the quadratic forms {$I_{1}$} at the boundaries and the quadratic forms {$I_{2}$} in the integral are positive definite. {Also,} there exists $\nu_{3}\in(0,\nu_{2})$
and $\delta_{4}\in(0,\delta_{3})$
such that, for any $\nu\in(0,\nu_{3})$,
and any $\delta\in(0,\delta_{4})$,
\begin{equation*}
\dot V\leq -\mu\min\limits_{[0,L]{\times[t_{0},+\infty)}}(\lambda_{1},\lambda_{2})V+C\delta V +C\bigl(\left(\lVert\boldsymbol{u}\rVert_{\infty}+\lVert\partial_{x}\boldsymbol{u}\rVert_{\infty}\right)^{{3}} \bigr),
\end{equation*}
where $C$ is a positive constant depending only on the system. Note that here, the cubic boundary terms that appeared in \eqref{dVtot} and the quadratic boundary terms proportional to $\delta$ have been compensated by the strictly negative quadratic boundary terms, taking $\delta$ and $\nu$ sufficiently small and using \eqref{estimate}. Note also that, thanks to \eqref{fluvial}, $\min_{[0,L]{\times[t_{0},+\infty)}}(\lambda_{1},\lambda_{2})>0$. Thus, choosing $\delta_{5}\in(0,\delta_{4})$ such that \hbox{$\delta_{5}<\mu\min_{[0,L]{\times[t_{0},+\infty)}}(\lambda_{1},\lambda_{2})/4C$}, for any $\delta\in(0,\delta_{5})$ one has
\begin{equation*}
\dot V\leq -\frac{3}{4}\,\mu\min\limits_{[0,L]{\times[t_{0},+\infty)}}(\lambda_{1},\lambda_{2})V+C \bigl(\left(\lVert\boldsymbol{u}\rVert_{\infty}+\lVert\partial_{x}\boldsymbol{u}\rVert_{\infty}\right)^{{3}} \bigr).
\end{equation*}
Now, if we assume in addition that \eqref{cV1} hold, using \eqref{estimate}, and Sobolev inequality, there exists $\nu_{4}\in(0,\nu_{3}]$ such that, for any $\nu\in(0,\nu_{4})$,
\begin{equation*}
C \bigl(\left(\lVert\boldsymbol{u}\rVert_{\infty}+\lVert\partial_{x}\boldsymbol{u}\rVert_{\infty}\right)^{{3}} \bigr)\leq \frac{\mu}{4}\min\limits_{[0,L]{\times[t_{0},+\infty)}}(\lambda_{1},\lambda_{2})V,
\end{equation*}
thus, setting $\gamma=\mu\min_{[0,L]{\times[t_{0},+\infty)}}(\lambda_{1},\lambda_{2})$,
\begin{equation*}
\dot V\leq -\frac{\gamma}{2}\, V{,\quad\forall t\in[t_{0},T]}
\end{equation*}
which shows the exponential decay of $V$ and \eqref{cV2}. {If in addition \eqref{cV1} holds, this}
ends the proof of Theorem \ref{th1}.

\subsection{Strategy to construct a Lyapunov function}
{A}ll that remains to do is to find $f_{1}$, $f_{2}$ and $q$ such that \eqref{c1}--\eqref{c3} are satisfied and such that $V$ satisfies \eqref{cV1}.
{To do so, we} first introduce the following function $\phi$ defined by
\begin{equation}
\begin{split}
\phi_{1}(t,x)&=\exp\biggl(\int_{0}^{x}\frac{\gamma_{1}}{\lambda_{1}}dx\biggr),\\
\phi_{2}(t,x)&=\exp\biggl(-\int_{0}^{x}\frac{\delta_{2}}{\lambda_{2}}dx\biggr),\\
\phi(t,x)&=\frac{\phi_{1}(t,x)}{\phi_{2}(t,x)},
\end{split}
\label{phi}
\end{equation}
The functions $\phi_{1}$ and $\phi_{2}$ are defined such that the diagonal change of variables $(y_{1} = \phi_{1}u_{1}$, $y_{2}= \phi_{2}u_{2})$ removes the diagonal source terms of the linearized system associated to \eqref{sys1} (recall that $\gamma_{1}$, $\delta_{2}$ are given in \eqref{l1l2}--\eqref{gamma} and are the diagonal coefficients of the source term of the linearized system). This change of variable is used for instance in \cite{BastinCoron22,C1_22} and is inspired from \cite[Chap.\,9]{KrsticBook}.\footnote{Although this is not used here, this change of variable also allows that the semigroup of the linearized system after change of variables has a compact difference with the semigroup of the homogeneous system, see \cite{hu2021minimal}.} In the following we are going to search for functions $f_{1}$, $f_{2}$ of the form
\begin{equation}
f_{1}(t,x)=\frac{\phi^{2}_{1}}{\lambda_{1}\xi(t,x)},\quad f_{2}(t,x)=\phi_{2}^{2}\,\frac{\xi(t,x)}{\lambda_{2}},
\label{deff1f2-0}
\end{equation}
where $\xi$ is a positive $C^{1}$ function to be defined. The motivation to look for functions of this form is the following: it was shown in \cite{BastinCoron22}, in the autonomous case with a purely proportional control, that \eqref{deff1f2-0} is an optimal choice in the following sense: if there exists a Lyapunov function of the form \eqref{def-V-tot}, then there exists a positive function $\xi\in C^{1}([0,L])$ such that \eqref{deff1f2-0} holds and the converse is true for suitable boundary conditions. This reduces the problem to finding a function $\xi$. Then, having~$f_{1}$ and~$f_{2}$ satisfying the differential inequalities \eqref{c3} is shown in this case (still in~\cite{BastinCoron22}) to be equivalent to having $\xi$ to be a supersolution of a given ODE. This is the approach followed for instance for the Saint-Venant equations in \cite{BC2017,HS} with a proportional control and in the autonomous case. Of course, the autonomous case with a proportional control is simpler than our current framework and brings some differences: the functions $\xi$, $f_{1}$ and $f_{2}$ do not depend on time. Also, because of this, we will want $\xi$ to satisfy
\begin{equation}
\label{eqxi}
\partial_{x}\xi >\Bigl|\frac{\phi\gamma_{2}}{\lambda_{1}}+\frac{\phi^{-1}\delta_{1}}{\lambda_{2}}\,\xi^{2}+{\frac{\phi}{\lambda_{1}^{2}}}\sqrt{\sfrac{g}{H_{1}}}\,\partial_{t} H_{1}\Bigr|,\quad \forall x\in[0,L],\ t\in[0,+\infty),
\end{equation}
and not anymore to be a supersolution of the ODE in \cite{BastinCoron22}.

In order to find such {a function $\xi$ and such functions $(f_{1}$, $f_{2})$, we start with} the following lemma, {from which we will construct a solution to \eqref{eqxi}. We will see later on that, if $\xi$ is a solution to \eqref{eqxi}, then $(f_{1}$, $f_{2})$ defined as \eqref{deff1f2-0} satisfy the condition~\eqref{c3}, at least under some condition on $\delta$.}
\begin{lem}
There exists $\delta_{0}>0$ such that if $\lVert\partial_{t}H_{1}\rVert_{L^{\infty}(({t_{0}},+\infty)\times(0,L)}\leq \delta_{0}$,
the function $\chi={\lambda_{2}\phi}/\lambda_{1}$ is solution on $[0,L]$ to the following equation
\begin{equation}
\partial_{x}\chi
=\Bigl|\frac{\phi\gamma_{2}}{\lambda_{1}}+\frac{\phi^{-1}\delta_{1}}{\lambda_{2}}\,\chi^{2}+{\frac{\phi}{\lambda_{1}^{2}}}\sqrt{\sfrac{g}{H_{1}}}\,\partial_{t} H_{1}\Bigr|, \quad\forall x\in[0,L],\ t\in[{t_{0}},+\infty),
\label{eqg}
\end{equation}
and for any $x\in[0,L]$ and any $t\in[{t_{0}},+\infty)$,
\begin{equation}
\Bigl(\frac{\phi\gamma_{2}}{\lambda_{1}}+\frac{\phi^{-1}\delta_{1}}{\lambda_{2}}\,\chi^{2}+{\frac{\phi}{\lambda_{1}^{2}}}\sqrt{\sfrac{g}{H_{1}}}\,\partial_{t} H_{1}\Bigr)>0.
\label{positivity}
\end{equation}
\label{lem2}
\end{lem}
The proof is given in Appendix {\ref{sec:prooflem2}}.
To understand the link between Lemma \ref{lem2} {and the choice of the coefficients of the Lyapunov function candidate $(f_{1}$, $f_{2})$, note that since $[0,L]$ is a closed interval we are going to be able to construct a solution $\chi_{\varepsilon}$ to \eqref{eqxi} from $\chi$. Then $(f_{1}$, $f_{2})$ will in turn be defined by \eqref{deff1f2-0}. This is what we do~now.}

\subsection{Construction of $\chi_{\varepsilon}$ and its properties}
Let now assume that $\delta \!<\!\delta_{0}$, where~$\delta_{0}$ is given by Lemma \ref{lem2}. {Recall that $\delta$ is the constant such that \eqref{boundHV} holds, namely
\begin{equation*}
\max\bigl(\lVert\partial_{t}H_{1}\rVert_{C^{1}([t_{0},\infty);{C^{0}([0,L])})},\lVert\partial_{t}V_{1}\rVert_{C^{1}([t_{0},\infty);{C^{0}([0,L])})}\bigr)<\delta.
\end{equation*}
As this is the only assumption on $(H_{1},V_{1})$, in the following we can assume without loss of generality that $t_{0}=0$.}
From Lemma \ref{lem2}, we know that there exists a solution on $[0,L]$ to equation \eqref{eqg}, which is ${\chi_{0}:=}\lambda_{2}\phi/\lambda_{1}$.
Therefore, as $[0,L]$ is a compact set, there exists $\varepsilon_{{0}}$ such that
for any $\varepsilon\in[0,\varepsilon_{{0}})$ there exists a solution $\chi_{\varepsilon}(t,x)$ to the following system
\begin{equation}
\begin{split}
\partial_{x}\chi_{\varepsilon}(t,x)&=\Bigl(\frac{\phi\gamma_{2}}{\lambda_{1}}+\frac{\delta_{1}}{\phi\lambda_{2}}(\chi_{\varepsilon})^{2}+{\frac{\phi}{\lambda_{1}^{2}}}\sqrt{\sfrac{g}{H_{1}}}\,\partial_{t} H_{1} \Bigr)+\varepsilon,\\
\chi_{\varepsilon}(0)&=\frac{\lambda_{2}(t,0)}{\lambda_{1}(t,0)}+\varepsilon,
\end{split}
\label{eqf}
\end{equation}
and moreover $(t,x,\varepsilon)\mto \chi_{\varepsilon}(t,x)$ is of class $C^{0}$ and $\partial_{x}\chi_{\varepsilon}(t,x)$ as well.
This is a classical result on {ODEs} due to Peano (see \eg \cite[Chap.\,5, Th.\,3.1]{Hartman}). Note that $\varepsilon_{0}>0$ a priori depends on $t$ and one could wonder whether $\varepsilon_{0}\to 0$ when $t\to+\infty$. We are going to show that this does not happen and we can choose $\varepsilon_{0}>0$ independent of $t\in[0,+\infty)$ such that $\chi_{\varepsilon}$ exists on $[0,+\infty)\times[0,L]$ for any $\varepsilon\in[0,\varepsilon_{0})$. Finally, note that $\chi_{\varepsilon}$ is a solution to \eqref{eqxi} if we can show that
\begin{equation*}
\Bigl(\frac{\phi\gamma_{2}}{\lambda_{1}}+\frac{\delta_{1}}{\phi\lambda_{2}}(\chi_{\varepsilon})^{2}+{\frac{\phi}{\lambda_{1}^{2}}}\sqrt{\sfrac{g}{H_{1}}}\,\partial_{t} H_{1} \Bigr)>0.
\end{equation*}
From \eqref{eqf},
$\partial_{t}\chi_{\varepsilon}$ satisfies the following equation
\begin{multline}
\label{eqdtf}
\partial_{x}\partial_{t}\chi_{\varepsilon}=2\frac{\delta_{1}}{\phi\lambda_{2}}\,\chi_{\varepsilon}\partial_{t}\chi_{\varepsilon}+\Bigl(\frac{\phi\gamma_{2}}{\lambda_{1}} \Bigr)_{t}+ \Bigl(\frac{\delta_{1}}{\phi\lambda_{2}} \Bigr)_{t}\chi^{2}_{\varepsilon}
+{\frac{\phi}{\lambda_{1}^{2}}}\sqrt{\sfrac{g}{H_{1}}}\,\partial_{tt}^{2} H_{1}\\
-{\frac{\phi}{\lambda_{1}^{2}}}\frac{1}{2}\sqrt{\sfrac{g}{H_{1}^{3}}}\,(\partial_{t}H_{1})^{2}+{\Bigl(\frac{\phi}{\lambda_{1}^{2}} \Bigr)_{t}\sqrt{\sfrac{g}{H_{1}}}\,\partial_{t} H_{1}}.
\end{multline}
We used here that, from Proposition \ref{propISS} and Remark \ref{rmkHpISS},
\[
(H_{1},V_{1})\in C^{0}([0,+\infty); H^{3}(0,L)),
\]
and from \eqref{target},
\begin{align*}
\partial_{tt}^{2}H_{1} &= -\partial_{tx}(H_{1}V_{1}), \quad\partial_{t}\partial_{x}H_{1}=-\partial_{x}^{2}(HV)\\
\tag*{and}
\partial_{t}\partial_{x}V_{1}&=\partial_{x}\left(-V_{1}\partial_{x}V_{1}-g\partial_{x}H_{1}-(kV_{1}^{2}/H_{1}-gC)\right).
\end{align*}
Thus $\partial_{tt}^{2}H_{1}$ belongs to $C^{0}([0,T]; H^{1}(0,L))$ and $(\gamma_{1}, \gamma_{2},\delta_{1},\delta_{2})$ to $C^{1}([0,T]; H^{1}(0,L))$. Besides, from \eqref{Lambda},
\begin{equation*}
\Bigl(\frac{\phi}{\lambda_{1}^{2}} \Bigr)_{t} = \frac{\partial_{t}\phi}{\lambda_{1}^{2}}-\frac{2\phi}{\lambda_{1}^{3}} \Bigl(\frac{{\sqrt{g}}\,\partial_{t}H_{1}}{2\sqrt{H_{1}}}+\partial_{t}V_{1} \Bigr).
\end{equation*}

Using \eqref{eqdtf}, we have
\begin{multline}
\label{difffeps}
\partial_{t}\chi_{\varepsilon}(t,x)=\partial_{t}\chi_{\varepsilon}(t,0)\exp\left(\int_{0}^{x}2\frac{\delta_{1}}{\phi\lambda_{2}}\,\chi_{\varepsilon}(t,y) dy\right)\\
\hspace*{-1cm}+\int_{0}^{x}\exp\left(\int_{y}^{x}2\frac{\delta_{1}}{\phi\lambda_{2}}\,\chi_{\varepsilon}(t,\omega) d\omega\right)\biggl(\Bigl(\frac{\phi\gamma_{2}}{\lambda_{1}}\Bigr)_{t}+\Bigl(\frac{\delta_{1}}{\phi\lambda_{2}}\Bigr)_{t}\chi^{2}_{\varepsilon}\\
+{\frac{\phi}{{\lambda_{1}^{2}}}}\sqrt{\sfrac{g}{H_{1}}}\,\partial_{tt}^{2} H_{1}-\frac{1}{2}{\frac{\phi}{{\lambda_{1}^{2}}}}\sqrt{\sfrac{g}{H_{1}^{3}}}\,(\partial_{t}H_{1})^{2}+{\Bigl(\frac{\phi}{\lambda_{1}^{2}}\Bigr)_{t}\sqrt{\sfrac{g}{H_{1}}}\,\partial_{t} H_{1}} \biggr)dy.
\end{multline}
Instead of seeing the function $\chi_{\varepsilon}$ as a solution of an ODE {(in space)} with a parameter~$t$, one can see it as a solution of an ODE with parameters
$\lambda_{1}$, $\lambda_{2}$, $\gamma_{2}$, $\delta_{1}$, $\partial_{t}H_{1}${, $\partial_{t}V_{1}$} and $\varepsilon$ that we denote ${\chi_{\varepsilon}(t,x) =} g_{\varepsilon}(x,\lambda_{1},\lambda_{2},\gamma_{1},\delta_{1},\partial_{t}H_{1}, {\partial_{t}V_{1}})$. From \cite[Th.\,2.1]{Hartman} the function
\begin{equation*}
(x,\varepsilon, \lambda_{1}, \lambda_{2}, \gamma_{2}, \delta_{1}, \partial_{t}H_{1}, {\partial_{t}V_{1}})\mto g_{\varepsilon}(x,\lambda_{1},\lambda_{2},\gamma_{1},\delta_{1},\partial_{t}H_{1}, {\partial_{t}V_{1}})
\end{equation*}
is continuous. But from \eqref{Hinfty}, \eqref{fluvial}, and \eqref{boundHV}, the quantities
\[
(\lambda_{1}(t), \lambda_{2}(t), \gamma_{2}(t), \delta_{1}(t), \partial_{t}H_{1}(t), {\partial_{t}V_{1}})
\]
are bounded from above and below and therefore belong to a compact set when $t\in[0,+\infty)$.
From this one can obtain that
\begin{equation}
\label{uniformconteps}
\varepsilon\mto g_{\varepsilon}(x,\lambda_{1}(t),\lambda_{2}(t),\gamma_{1}(t),\delta_{1}(t),\partial_{t}H_{1}(t), {\partial_{t}V_{1}})=\chi_{\varepsilon}(t,x)
\end{equation}
is uniformly continuous in $\varepsilon \;{\in [0,\varepsilon_{1})}$ for $(t,x)\in[0,{+}\infty)\times[0,L]$, for some $\varepsilon_{1}>0$ (independent of $t$). One can obtain this uniform continuity and this $\varepsilon_{1}$ by observing that, from Lemma \ref{lem2} and \eqref{eqf},
\begin{equation*}
\begin{split}
\partial_{x}(\chi_{\varepsilon}(t,x)-\chi_{0}(t,x)) &= \frac{\delta_{1}}{\phi\lambda_{2}}(\chi_{\varepsilon}^{2}-\chi_{0}^{2})+\varepsilon=\frac{\delta_{1}}{\phi\lambda_{2}}\left[(\chi_{\varepsilon}-\chi_{0})^{2}+2\chi_{0}(\chi_{\varepsilon}-\chi_{0})\right]+\varepsilon\\
&\leq C\left[(\chi_{\varepsilon}-\chi_{0})^{2}+(\chi_{\varepsilon}-\chi_{0})\right]+\varepsilon,
\end{split}
\end{equation*}
where $C$ is a positive constant that might change between lines but depends only on $H_{\max}$, $\alpha$ and an upper bound of $\delta$ and where we used that $\chi_{0} = \lambda_{2}\phi/\lambda_{1}$. Let us set $h$ defined by
\begin{equation}
\label{defhgr}
\begin{split}
\partial_{x}h(t,x) &= C\left[h^{2}(t,x)+h(t,x)\right]+\varepsilon,\\
h(t,0) &= \chi_{\varepsilon}(t,0)-\chi_{0}(t,0) = \varepsilon,
\end{split}
\end{equation}
{by comparison} $\chi_{\varepsilon}(t,x)-\chi_{0}(t,x)\leq h(t,x)$ for any $(t,x)\in[0,+\infty)\times [0,L]$, provided~$h$ exist on $[0,+\infty)\times [0,L]$. This implies that $\chi_{\varepsilon}$ exist on $[0,+\infty)\times [0,L]$ provided $h$ does. Besides, from \eqref{defhgr}, $h$ is nondecreasing and therefore $h(t,x)\geq h(t,0)=\varepsilon$ so
\begin{equation*}
\partial_{x}h(t,x) \leq 2C\bigl[h^{2}(t,x)+h(t,x)\bigr],
\end{equation*}
which implies, integrating and using that $\ln(x/1+x)$ is a primitive of $1/(x^{2}+x)$,
\begin{equation}
\label{exprheq}
\frac{h(t,x)}{1+h(t,x)}\leq \frac{\varepsilon}{1+\varepsilon}\,e^{2CL},
\end{equation}
From \eqref{exprheq}
we deduce that there exists
$\varepsilon_{1}$ depending only on $H_{\max}$, $\alpha$, $L$ and an upper bound of $\delta$ such that for any $\varepsilon\in[0,\varepsilon_{1})$, $h$ (and hence $\chi_{\varepsilon}$) exists on $[0,+\infty)\times[0,L]$ and
\begin{equation*}
\chi_{\varepsilon}(t,x)-\chi_{0}(t,x)\leq h(t,x)\leq C(\varepsilon),\quad\forall (t,x)\in[0,+\infty)\times[0,L],
\end{equation*}
where $C(\varepsilon)$ depends only on $H_{\max}$, $\alpha$, $L$, $\varepsilon$, and an upper bound of $\delta$, varying continuously with $\varepsilon$ with $C(0)=0$.
By {comparison} again and \eqref{eqf} we have $(\chi_{\varepsilon}-\chi_{0})\geq 0$, hence
\begin{equation*}
|\chi_{\varepsilon}(t,x)|\leq |\chi_{0}(t,x)|+ C(\varepsilon),\quad\forall (t,x)\in[0,+\infty)\times[0,L].
\end{equation*}
Finally we conclude to the uniform continuity of $\varepsilon\mto \chi_{\varepsilon}$ with respect $(t,x)\in[0,+\infty)\times[0,L]$ for $\varepsilon \in [0,\varepsilon_{1})$ by using the fact that $\chi_{0}=\lambda_{2}\phi/\lambda_{1}$ is uniformly bounded on $[0,+\infty)\times[0,L]$. Indeed, from \eqref{fluvial} and \eqref{Hinfty} we know that for any $(t,x)\in[0,{+}\infty)\times[0,L]$,
\begin{equation}
\sqrt{g H_{\max}}>\lambda_{2}>\alpha,\quad2\sqrt{gH_{\max}}>\lambda_{1}>\alpha.
\label{bornel1l2}
\end{equation}
Besides, from the definition of $\phi_{1}$ and $\phi_{2}$ given by \eqref{phi}, \eqref{gamma} and the bound \eqref{fluvial}, \eqref{Hinfty}, there exists a constant $C_{8}$ that only depends on
$\delta$, $\alpha$ and $H_{\max}$ such that
\begin{equation}
\frac{1}{C_{8}}\leq\lVert\phi_{1}\rVert_{\infty}\leq C_{8},\quad\frac{1}{C_{8}}\leq\lVert\phi_{2}\rVert_{\infty}\leq C_{8}.
\label{bornephi1phi2}
\end{equation}
This, together with \eqref{difffeps} implies that there exists $C_{0}$ depending only on $L$, $H_{\max}$, $\alpha$, $\varepsilon$, {an upper bound of $\delta$ (for instance $\delta_{0}$),} and continuous with $\varepsilon\in[0,\varepsilon_{1})$ such that
\begin{multline*}
\left|\int_{0}^{x}\exp\left(\int_{y}^{x}2\frac{\delta_{1}}{\phi\lambda_{2}}\,\chi_{\varepsilon}(t,\omega) d\omega\right)\partial_{t}(\partial_{y}(H_{1}V_{1}))dy\right|\\
\leq C_{0} \max \bigl(\lVert\partial_{t}H_{1}\rVert_{C^{1}([0,+\infty);C^{0}([0,L]))},\lVert\partial_{t}V_{1}\rVert_{C^{1}([0,+\infty);C^{0}([0,L]))}\bigr).
\end{multline*}
Similarly
there exists a constant $C_{1}>0$ depending only on $L$, $H_{\max}$, $\alpha${, and an upper bound of $\delta$} such that
\begin{multline}
\label{dphi}
\rVert \partial_{t}{\phi}\lVert_{L^{\infty}((0,+\infty)\times(0,L)}\\
\leq C_{1} \max \bigl(\lVert\partial_{t}H_{1}\rVert_{C^{1}([0,+\infty);C^{0}([0,L]))},\lVert\partial_{t}V_{1}\rVert_{C^{1}([0,+\infty);C^{0}([0,L]))} \bigr),
\end{multline}
and similarly for $\phi_{2}$. This, together with the definition of $\lambda_{1}$ and $\lambda_{2}$ given by \eqref{Lambda}, \eqref{difffeps}, and
using the continuity of $\varepsilon\to \chi_{\varepsilon}$ on $[0,\varepsilon_{1})$ (recall that this continuity is uniform with respect to $(t,x)\in[0,+\infty)\times[0,L]$ {from \eqref{uniformconteps}}),
we get that there exists $C>0$ depending only on $H_{\max}$, $\alpha$, {an upper bound of $\delta$,} $\varepsilon$ and continuous with $\varepsilon$ on $[0,\varepsilon_{1})$ such that\vspace*{-5pt}
\begin{multline*}
|\partial_{t}\chi_{\varepsilon}(t,x)|\leq \Bigl(|\partial_{t}\chi_{\varepsilon}(t,0)|\\
+ \max \bigl(\lVert\partial_{t}H_{1}\rVert_{C^{1}([0,+\infty);C^{0}([0,L]))},\lVert\partial_{t}V_{1}\rVert_{C^{1}([0,+\infty);C^{0}([0,L]))} \bigr) \Bigr)C(\varepsilon).
\end{multline*}
But, from \eqref{eqf} $\partial_{t}\chi_{\varepsilon}(t,0)=\left(\lambda_{2}{(0)}/\lambda_{1}{(0)}\right)_{t}$, thus
using \eqref{boundHV}
we obtain
\begin{equation}
|\partial_{t}\chi_{\varepsilon}(t,x)|\leq \delta C_{2}(\varepsilon),
\label{borne_dtfeps}
\end{equation}
where $C_{2}$ is again a constant that only depends on $\varepsilon$, $\alpha$, $H_{\max}$, {an upper bound of $\delta$} and is continuous with $\varepsilon$ on $[0,\varepsilon_{1})$.
We can now restrict ourselves to $\varepsilon\in[0,\varepsilon_{1}/2]$ and then $C_{2}$ can be chosen independent of $\varepsilon$ by simply taking its maximum on $[0,\varepsilon_{1}/2]$.
Recall that from Lemma \ref{lem2} we have, $\chi_{0}=\phi\lambda_{2}/\lambda_{1}$, and
\begin{equation*}
\Bigl(\frac{\phi\gamma_{2}}{\lambda_{1}}+\frac{\delta_{1}}{\phi\lambda_{2}}\,\chi_{0}^{2}+{\frac{\phi}{\lambda_{1}^{2}}}\sqrt{\sfrac{g}{H_{1}}}\,\partial_{t} H_{1} \Bigr)>0.
\end{equation*}
Recall {also} that we did not yet choose the bound $\delta\in(0,{\delta_{0}})$ on $\lVert\partial_{t}H_{1}\rVert_{C^{1}([0,\infty);C^{0}([0,L]))}$ and $\lVert\partial_{t}V_{1}\rVert_{C^{1}([0,\infty);C^{0}([0,L]))}$ given in \eqref{boundHV}.
From the assumptions on $k_{p}$ and $k_{I}$, \ie \eqref{cond}, and \eqref{k1k3}, {and recalling that $k_{1}=\partial_{1}\mathcal{D}_{2}(0,0,t)$
and $k_{3}=-\partial_{2}\mathcal{D}_{2}(0,0,t)$}, one has
\begin{equation}
k_{1}^{2}<\Bigl(\frac{\lambda_{1}(L)}{\lambda_{2}(L)} \Bigr)^{2},\quad k_{3}>0.
\label{condk1k3}
\end{equation}
Thus, using \eqref{Lambda},
\begin{equation}
\label{defeta1}
\eta_{1}:=\min \biggl(\Bigl(\frac{1}{|k_{1}|}-\frac{\lambda_{2}(L)}{\lambda_{1}(L)} \Bigr),1-\frac{\lambda_{2}(L)}{\lambda_{1}(L)} \biggr)>0.
\end{equation}
As $\varepsilon\to \chi_{\varepsilon}(t,x)$ is uniformly continuous with $\varepsilon$ for $(t,x)\in[0,{+}\infty)\times[0,L]$,
there exists
$\varepsilon_{2}\in(0,\varepsilon_{1}/2)$ such that
for any $(t,x)\in[0,{+}\infty)\times[0,L]$
\begin{equation}
|\chi_{\varepsilon_{2}}(t,x)-\chi_{0}(t,x)| \;{<}\; \phi(t,L)\eta_{1},
\label{fL}
\end{equation}
and
\begin{equation}
\Bigl(\frac{\phi\gamma_{2}}{\lambda_{1}}+\frac{\delta_{1}}{\phi\lambda_{2}}\,\chi_{\varepsilon_{2}}^{2}+{\frac{\phi}{\lambda_{1}^{2}}}\sqrt{\sfrac{g}{H_{1}}}\,\partial_{t} H_{1} \Bigr)>0.
\label{fpositive}
\end{equation}
{In particular $\chi_{\varepsilon}$ is a solution to \eqref{eqxi}.}
Note that $\varepsilon_{2}$ depends a priori on $\delta$ from \eqref{fpositive}. However, from Lemma \ref{lem2}
we can in fact choose $\varepsilon_{2}$ independent of $\delta$ and depending only on an upper bound of $\delta$ (for instance $\delta_{0}$ given by Lemma \ref{lem2}). This is important as, in the following,
we will choose
a $\delta$ that may depends on $\varepsilon$.

\subsection{Condition from the integral}
{As announced we} select $f_{1}$ and $f_{2}$ in the following way:
\begin{equation}
\begin{split}
f_{1}(t,x)&=\frac{\phi^{2}_{1}}{\lambda_{1}\chi_{\varepsilon_{2}}(t,x)}>0,\\
f_{2}(t,x)&=\phi_{2}^{2}\,\frac{\chi_{\varepsilon_{2}}(t,x)}{\lambda_{2}}>0,
\label{deff1f2}
\end{split}
\end{equation}
and we can now check that the condition \eqref{c3} is verified for $\delta$ small enough. {We have}
\begin{equation*}
(-\lambda_{1}f_{1})_{x}=-2\,\frac{(\phi_{1})_{x}\lambda_{1}f_{1}}{\phi_{1}}+\phi_{1}^{2}\,\frac{\partial_{x}\chi_{\varepsilon_{2}}(t,x)}{\chi_{\varepsilon_{2}}^{2}(t,x)}.
\end{equation*}
Thus from \eqref{phi}
\begin{equation}
-(\lambda_{1}f_{1})_{x}+2\gamma_{1}f_{1}
=\phi_{1}^{2}\frac{\partial_{x}\chi_{\varepsilon_{2}}}{\chi^{2}_{\varepsilon_{2}}}
\label{f1}
\end{equation}
and similarly
\begin{equation}
\begin{split}
(\lambda_{2}f_{2})_{x}+2\delta_{2}f_{2}&=(\phi^{2}_{2}\chi_{\varepsilon_{2}}(t,x))_{x}-(\phi^{2}_{2})_{x}\chi_{\varepsilon_{2}}(t,x)\\
&=\phi_{2}^{2}\partial_{x}\chi_{\varepsilon_{2}}.
\end{split}
\label{f2}
\end{equation}
Therefore, from \eqref{eqf}, \eqref{f1}, and \eqref{f2}, one has
\begin{multline}
\label{cond00}
(-(\lambda_{1}f_{1})_{x}+2\gamma_{1}f_{1}-\partial_{t}f_{1})((\lambda_{2}f_{2})_{x}+2\delta_{2}f_{2}-\partial_{t}f_{2})\\
=\Bigl(\frac{\phi_{1}\phi_{2}}{\chi_{\varepsilon_{2}}} \Bigr)^{2}\biggl(\Bigl(\frac{\phi \gamma_{2}}{\lambda_{1}}+\frac{\delta_{1}}{\phi\lambda_{2}}\,\chi_{\varepsilon_{2}}^{2}+{\frac{\phi}{\lambda_{1}^{2}}}\sqrt{\sfrac{g}{H_{1}}}\,\partial_{t} H_{1} \Bigr)
+\varepsilon_{2} \biggr)^{2}\\
-\partial_{x}\chi_{\varepsilon_{2}} \Bigl(\frac{\phi_{1}^{2}}{\chi^{2}_{\varepsilon_{2}}}\,\partial_{t}f_{2}
+\phi_{2}^{2}\partial_{t}f_{1} \Bigr)
+(\partial_{t}f_{1})(\partial_{t}f_{2}).
\end{multline}
But we have
\begin{equation}
\label{dtf1}
\partial_{t}f_{1}=2\frac{(\partial_{t}\phi_{1})\phi_{1}}{\lambda_{1}\chi_{\varepsilon_{2}}}-\Bigl(\frac{\partial_{t}\lambda_{1}}{\lambda_{1}^{2}\chi_{\varepsilon_{2}}}+\frac{\partial_{t}\chi_{\varepsilon_{2}}}{\lambda_{1}\chi^{2}_{\varepsilon_{2}}} \Bigr)\phi_{1}^2,
\end{equation}
and besides, from \eqref{target} and \eqref{boundHV}, there exists $C_{3}>0$ depending only on $\alpha$ and $H_{\max}$, and an upper bound of $\delta$ (for instance ${\delta_{0}}$), such that
\begin{equation*}
\max\bigl(\lVert H_{1x}\rVert_{L^{\infty}((0,+\infty)\times(0,L)},\lVert V_{1x}\rVert_{L^{\infty}((0,+\infty)\times(0,L)}\bigr)\leq C_{3}.
\end{equation*}
Thus, using \eqref{gamma} and \eqref{Lambda}, there exists $C_{4}>0$ depending only on {$L$,} $\alpha$ and $H_{\max}$, and $\delta_{0}$ (but not on $\delta$) such that
\begin{equation}
\max \bigl(\lVert\phi_{1}\rVert_{L^{\infty}((0,+\infty)\times(0,L){)}},\lVert\phi^{-1}_{{1}}\rVert_{L^{\infty}((0,+\infty)\times(0,L){)}} \bigr)<C_{4},
\label{boundphi}
\end{equation}
and similarly for $\phi_{2}$.
Observe now that, from $\chi_{0}=\lambda_{2}\phi/\lambda_{1}$ and \eqref{boundphi}, $|\chi_{0}|$ and $1/|\chi_{0}|$ can be bounded by a constant depending only on {$L$,} $\alpha$, $H_{\max}$, and
$\delta_{0}$. Thus from
\eqref{fL} one obtains
\begin{equation}
1/C_{5}\leq \lVert \chi_{\varepsilon_{2}}\rVert_{L^{\infty}((0,+\infty)\times(0,L)}\leq C_{5},
\label{boundf}
\end{equation}
where $C_{5}$ only depends on {$L$,} $\alpha$, $H_{\max}$ and $\delta_{0}$.
{Therefore}, from
\eqref{Lambda}, \eqref{borne_dtfeps}, \eqref{dphi}, {\eqref{dtf1}} and \eqref{boundf} one has
\begin{equation*}
|\partial_{t}f_{1}|\leq C_{6}\delta,
\end{equation*}
and similarly
\begin{equation*}
|\partial_{t}f_{2}|\leq C_{7}\delta,
\end{equation*}
where $C_{6}$ and $C_{7}$ are constants that only depend on {$L$,} $\alpha$, $H_{\max}$ (and $\delta_{0}$).
We now select the bound on
$\max \bigl({\lVert\partial_{t}H_{1}\rVert_{C^{1}([t_{0},\infty);{C^{0}([0,L])})},\lVert\partial_{t}V_{1}\rVert_{C^{1}([t_{0},\infty);{C^{0}([0,L])})}} \bigr)$:
we select $\delta_{3}\in(0,\delta_{0})$ such that, for any $\delta\in[0,\delta_{3}]$ and any $(t,x)\in[0,{+}\infty)\times[0,L]$,
\begin{equation}
\label{conddelta02}
C_{6}C_{5}^{2}C_{4}^{2}\delta<\varepsilon_{2},
\end{equation}
and
\begin{multline}
\label{conddelta2}
\varepsilon_{2}^{2}+2\varepsilon_{2}\inf_{\substack{x\in[0,L]\\ t\in[0,+\infty)\\\varepsilon\in(0,\varepsilon_{2})}}\Bigl(\frac{\phi \gamma_{2}}{\lambda_{1}}+\frac{\delta_{1}}{\phi\lambda_{2}}\chi_{\varepsilon}^{2}+{\frac{\phi}{\lambda_{1}^{2}}}\sqrt{\sfrac{g}{H_{1}}}\,\partial_{t} H_{1} \Bigl)\\
> \Bigl(\frac{\phi \gamma_{2}}{\lambda_{1}}+\frac{\delta_{1}}{\phi\lambda_{2}}\,X^{2}+{\frac{\phi}{\lambda_{1}^{2}}}\sqrt{\sfrac{g}{H_{1}}}\,\delta+\varepsilon_{2} \Bigr) \Bigl(C_{6}\frac{\phi_{1}^{2}}{X^{2}}+C_{7}\phi_{2}^{2} \Bigr) \Bigl(\frac{X}{\phi_{1}\phi_{2}} \Bigr)^{2}\delta\\
+2\,\frac{\phi}{\lambda_{1}^{2}}\sqrt{\sfrac{g}{H_{1}}} \Bigl(\frac{\phi\gamma_{2}}{\lambda_{1}}+\frac{\phi^{-1}\delta_{1}}{\lambda_{2}}\,X^{2} \Bigr)\delta+\biggl(\Bigl(\frac{X}{\phi_{1}\phi_{2}} \Bigr)^{2}C_{7}C_{6}{+\frac{\phi^{2}}{\lambda_{1}^{4}}\frac{g}{H_{1}}}\biggr)\delta^{2},
\end{multline}
for any $x\in[0,L]$ and any $X\in[1/C_{5},C_{5}]$ ({note} that having it for $X=C_{5}$ is enough). This is possible as $\varepsilon_{2}>0$ and, when $\delta_{3}=0$,
\eqref{conddelta2} is verified and the inequality is strict.
Then, from \eqref{phi}, \eqref{boundphi}, \eqref{cond00}, \eqref{boundf}--\eqref{conddelta2},
\begin{equation*}
\begin{split}
(-(\lambda_{1}f_{1})_{x}+2\gamma_{1}f_{1}-\partial_{t}f_{1})((\lambda_{2}f_{2})_{x}+2\delta_{2}f_{2}-\partial_{t}f_{2})&>
\Bigl(\frac{\phi_{1}\phi_{2}}{\chi_{\varepsilon_{2}}} \Bigr)^{2} \Bigl(\frac{\phi \gamma_{2}}{\lambda_{1}}+\frac{\delta_{1}}{\phi\lambda_{2}}\,\chi_{\varepsilon_{2}}^{2} \Bigr)^{2}\\
&={\left(\gamma_{2}f_{1}+\delta_{1}f_{2}\right)^{2}},\\
\end{split}
\end{equation*}
which is exactly the second inequality of \eqref{c3}. Besides,
from \eqref{positivity} and \eqref{conddelta02},
\begin{multline*}
(-(\lambda_{1}f_{1})_{x}+2\gamma_{1}f_{1}-\partial_{t}f_{1})=\phi_{1}^{2}\frac{\partial_{x}\chi_{\varepsilon_{2}}}{\chi^{2}_{\varepsilon_{2}}}
-\partial_{t}f_{1}
\\
=\frac{\phi_{1}^{2}}{\chi_{\varepsilon_{2}}^{2}}\biggl(
\Bigl(\frac{\phi\gamma_{2}}{\lambda_{1}}+\frac{\delta_{1}}{\phi\lambda_{2}}\,{\chi}_{{\varepsilon_{2}}}^{2}+{\frac{\phi}{\lambda_{1}^{2}}}\sqrt{\sfrac{g}{H_{1}}}\,\partial_{t} H_{1} \Bigl)+{\varepsilon_{2}}
-\frac{\partial_{t}f_{1}\chi_{\varepsilon_{2}}^{2}}{\phi_{1}^{2}}
\biggl)>0.
\end{multline*}

\subsection{Conditions at the boundaries}
We can now check that
\eqref{c1} and \eqref{c2} are also verified thanks to the choice of $\varepsilon_{2}$ and $\eta_{1}${, and $(f_{1}$, $f_{2})$ given by \eqref{deff1f2}}. Indeed, using \eqref{eqf}{, \eqref{phi}} and \eqref{k1k3}, one has
\begin{equation*}
\frac{\lambda_{2}(0)f_{2}(t,0)}{\lambda_{1}(0)f_{1}(t,0)}=\chi_{\varepsilon_{2}}^{2}(t,0)= \Bigl(\frac{\lambda_{2}(0)}{\lambda_{1}(0)}+\varepsilon_{2} \Bigr)^{2}> \Bigl(\frac{\lambda_{2}(0)}{\lambda_{1}(0)} \Bigr)^{2}=k_{2}^{2},
\end{equation*}
{which is exactly \eqref{c1}.} This explains our choice of initial condition for $\chi_{\varepsilon_{2}}$.
Now, from \eqref{fL}, one has
\begin{equation*}
\frac{\lambda_{1}(t,L)f_{1}(t,L)}{\lambda_{2}(t,L)f_{2}(t,L)}=\frac{\phi^{2}(t,L)}{\chi^{2}_{\varepsilon_{2}}(L)}>\frac{1}{\left(\sfrac{\lambda_{2}(t,L)}{\lambda_{1}(t,L)}+\eta_{1}\right)^{2}},
\end{equation*}
and from the definition of $\eta_{1}$ given by \eqref{defeta1},
\begin{equation*}
\eta_{1}+\frac{\lambda_{2}(L)}{\lambda_{1}(L)}=\min(\sfrac{1}{|k_{1}|},1).
\end{equation*}
Therefore,
\begin{equation}
\frac{\lambda_{1}(t,L)f_{1}(t,L)}{\lambda_{2}(t,L)f_{2}(t,L)}>\max(k_{1}^{2},1),
\label{cond3}
\end{equation}
and in particular
the condition \eqref{c2a} is verified. Let us now look at condition \eqref{c2b}. So far we have not selected the positive constant $q$.
We want to show that there exists $q>0$ such that the condition \eqref{c2b} is satisfied. Observe that the left-hand side of \eqref{c2b} can be seen as a polynomial in $q$,
and the condition \eqref{c2b} can be rewritten as
\begin{equation*}
\begin{split}
P(q)&:=
-\frac{q^{2}}{4}\frac{H_{1}}{g}\left(k_{1}-1\right)^{2}+q\sqrt{\sfrac{H_{1}}{g}}\,k_{3}\left(\lambda_{1}f_{1}(L)-\lambda_{2}f_{2}(L)(k_{1}^{2}-k_{1}(k_{1}-1))\right)
\\
&\hspace*{3.23cm}-\left(\lambda_{1}f_{1}(L)\right)\left(\lambda_{2}f_{2}(L)\right)k_{3}^{2}\\
&\,=-\frac{q^{2}}{4}\frac{H_{1}}{g}\left(k_{1}-1\right)^{2}+q\sqrt{\sfrac{H_{1}}{g}}\,k_{3}\left(\lambda_{1}f_{1}(L)-\lambda_{2}f_{2}(L)k_{1})\right)\\
&\hspace*{3.23cm}-\left(\lambda_{1}f_{1}(L)\right)\left(\lambda_{2}f_{2}(L)\right)k_{3}^{2}\\
&\,>0.
\end{split}
\end{equation*}
From \eqref{cond3} $\lambda_{1}f_{1}(t,L)>\lambda_{2}f_{2}(t,L)k_{1}$ and from \eqref{condk1k3} $k_{3}>0$. Thus the real roots of~$P$ are positive if they exist. This implies that there exists a positive constant $q$ such that \eqref{c2b} is satisfied if the discriminant of $P$
is positive. Denoting its discriminant by $\Delta$,
\begin{equation*}
\Delta=\frac{H_{1}}{g}\,k_{3}^{2}\lambda_{2}^{2}f_{2}^{2}(t,L)\biggl[\Bigl(\frac{\lambda_{1}f_{1}(L)}{\lambda_{2}f_{2}(L)}-k_{1} \Bigr)^{2}-\Bigl(\frac{\lambda_{1}f_{1}(L)}{\lambda_{2}f_{2}(L)} \Bigr)\left(k_{1}-1\right)^{2}\biggr].
\end{equation*}
Let us introduce $h:X\to (X-k_{1})^{2}-X(k_{1}-1)^{2}$. The function $h$ is a second order polynomial with a positive dominant coefficient and observe that its roots are
$k_{1}^{2}$ and~$1$. Thus
$h$ is increasing strictly on $[\max(k_{1}^{2},1),+\infty)$.
Hence, using \eqref{cond3},
\begin{equation*}
\begin{split}
\Delta&=\frac{H_{1}}{g}\,k_{3}^{2}\lambda_{2}^{2}f_{2}^{2}(t,L)h \Bigl(\frac{\lambda_{1}f_{1}(L)}{\lambda_{2}f_{2}(L)} \Bigr)\\
&>\frac{H_{1}}{g}\,k_{3}^{2}\lambda_{2}^{2}f_{2}^{2}(t,L)h(\max(k_{1}^{2},1))
=0.
\end{split}
\end{equation*}
This proves that there exists $q>0$ such that \eqref{c2b} is satisfied, and we select such $q$.
All it remains to do now is to show that the function $(\boldsymbol{U},z,{t})\to V(\boldsymbol{U},z,{t})$, which is now entirely selected, satisfies \eqref{cV1}.

Thus, using that $\chi_{0}=\lambda_{2}\phi/\lambda_{1}$, \eqref{deff1f2}, \eqref{fL}, \eqref{bornephi1phi2}, and \eqref{bornel1l2}, there exists $\eta>0$, $c_{1}>0$ constant independent of $\boldsymbol{U}$ and ${Z}$ such that, for any $(\boldsymbol{U},{Z})\in H^{2}(0,L)\times \mathbb{R}$ {with $\|\boldsymbol{U}\|_{H^{2}}+|Z|\leq \eta$}
\begin{equation*}
c_{1}\left(\lVert \boldsymbol{U}\rVert_{H^{2}(0,L)}+|Z|\right) \leq V(\boldsymbol{U},Z,{t})\leq \frac{1}{c_{1}}\left(\lVert \boldsymbol{U}\rVert_{H^{2}(0,L)}+|Z|\right)\quad\forall t\in[0,+\infty),
\end{equation*}
which is exactly \eqref{cV1}. This concludes the proof of Theorem \ref{th1}.
\section{Conclusion}
In this paper, we gave simple conditions on the design of a single PI controller to ensure
the exponential stability of the nonlinear Saint-Venant equations with arbitrary friction and slope in the $H^{2}$ norm.
These conditions apply when the inflow is an unknown constant{. I}n that case the system has steady-states and any of them is stable. {Additionally these conditions} also apply when
the inflow is time-dependent and slowly variable (with potentially a large total variation). In that case, no steady-state exists and one has to stabilize other target states.
When the values of the target state are known at the end of the river, we have exponential stability of the target state.
{In other situations}, we have the Input-to-State Stability with respect to the variation of the inflow disturbance.
These sufficient conditions are found using a local quadratic entropy
and, to the best of our knowledge, are less restrictive than any of the conditions that existed so far, even in the linear case.
In \cite{BastinCoronTamasoiu2015} it was shown that, in absence of friction and slope, these conditions {are} optimal for the linear case. When there is some slope or friction, however, there is so far no answer{. Knowing }whether the conditions of Theorem \ref{th1} are optimal or not would be a very interesting {open question} for a further study. {Its possible application to a network of channels would also be a matter of interest.
Finally, many stabilizing devices for finite dimensional systems also use a PID control with an additional derivative term.
It has been shown in \cite{CoronTamasoiu2015} that this control cannot ensure exponential stability for a homogeneous hyperbolic equation.
It would be an interesting question to know whether a filtering on the derivative term could enable to recover the stability for infinite dimensional system and whether this would enable a faster stabilization than the PI control.}

\appendix
\section{Boundary conditions (\ref{bound}) and (\ref{k1k3})}
\label{boundary}
In this appendix we justify the boundary conditions \eqref{bound} with \eqref{k1k3} after the change of variables. From the boundary conditions \eqref{bound1} in the physical coordinate $(h,v)$, together with the definition of $u_{1}$ and $u_{2}$
given in \eqref{change2}, one has at $x=L$
\begin{equation}
\begin{split}
u_{1}(t,L)=\mathcal{B}_{2}(h(t,L),Z(t),t)+\sqrt{\sfrac{g}{H_{1}}}\,h(t,L)=:\mathcal{F}_{1}(h(t,L),Z(t),x,t),\\
u_{2}(t,L)=\mathcal{B}_{2}(h(t,L),Z(t),t)-\sqrt{\sfrac{g}{H_{1}}}\,h(t,L)=:\mathcal{F}_{2}(h(t,L),Z(t),x,t).
\end{split}
\label{Bcond1}
\end{equation}
From its definition, $\mathcal{F}_{1}$ is $C^{1}$ and, from \eqref{k1k30}, and \eqref{estimate}, there exists $\nu_{1}\in(0,\nu_{0})$ such that, for any $t\in[0,\infty)$, $\partial_{1}{\mathcal{F}_{1}}(0,Z(t),t)\neq 0$. Thus $\mathcal{F}_{1}$ is locally invertible with respect to its first variable, thus
there exists $\nu_{2}\in(0,{\nu_{1}})$ such that
$h(t,L)=\mathcal{F}_{1}^{-1}(u_{1}(t,L),Z(t),t)$, where $\mathcal{F}_{1}^{-1}$ denotes the inverse with respect to the first variable. Besides, as $\mathcal{F}_{1}$ is of class $C^{2}$ with respect to the two first variables, $\mathcal{F}_{1}^{-1}$ is also of class $C^{2}$. Then, using \eqref{Bcond1}
\begin{equation*}
u_{2}(t,L)=\mathcal{F}_{2}(\mathcal{F}_{1}^{-1}(u_{1}(t,L),Z(t),t),Z(t),t)=:\mathcal{D}_{2}(u_{1}(t,L),Z(t),t),
\end{equation*}
and, using \eqref{k1k30},
\begin{equation*}
\begin{split}
\partial_{1}\mathcal{D}_{2}(0,0,t)&=\partial_{1}\mathcal{F}_{2}(0,0,t)\partial_{1}(\mathcal{F}_{1}^{-1})(0,0,t)\\
&=\frac{\partial_{1}\mathcal{F}_{2}(0,0,t)}{\partial_{1}\mathcal{F}_{1}(0,0,t)}=\frac{\partial_{1}\mathcal{B}_{2}(0,0,t)-\sqrt{\sfrac{g}{H_{1}}}}{\partial_{1}\mathcal{B}_{2}(0,0,t)+\sqrt{\sfrac{g}{H_{1}}}}=-\frac{\lambda_{1}(L)-v_{G}(1+k_{p})}{\lambda_{2}(L)+v_{G}(1+k_{p})}.
\end{split}
\end{equation*}
Now, as $\partial_{2}\mathcal{F}_{1}^{-1}(0,0,t)=-\partial_{2}\mathcal{F}_{1}(0,0,t)/\partial_{1}\mathcal{F}_{1}(0,0,t)$, using \eqref{k1k30},
\begin{equation*}
\begin{split}
\partial_{2}\mathcal{D}_{2}(0,0,t)&=\partial_{1}\mathcal{F}_{2}(0,0,t)\partial_{2}(\mathcal{F}_{1}^{-1})(0,0,t)+\partial_{2}\mathcal{F}_{2}(0,0,t)\\
&=-\partial_{1}\mathcal{F}_{2}(0,0,t)\frac{\partial_{2}\mathcal{F}_{1}(0,0,t)}{\partial_{1}\mathcal{F}_{1}(0,0,t)}+\partial_{2}\mathcal{F}_{2}(0,0,t)\\
&=\partial_{2}\mathcal{B}_{2}(0,0,t)\Bigl(1-\frac{\partial_{1}\mathcal{B}_{2}(0,0,t)-\sqrt{\sfrac{g}{H_{1}}}}{\partial_{1}\mathcal{B}_{2}(0,0,t)+\sqrt{\sfrac{g}{H_{1}}}}\Bigr)\\
&=-\frac{v_{G}k_{I}}{H_{1}(t,L)} \Bigl(\frac{2\sqrt{gH_{1}(t,L)}}{v_{G}(1+k_{p})+\lambda_{2}(t,L)} \Bigr).
\end{split}
\end{equation*}
The same can be done in $x=0$ in a slightly easier way, as $\mathcal{B}_{1}$ does not depends on~$Z$. This gives \eqref{bound} and \eqref{k1k3}.

\section{Proof of Proposition \ref{ISS}}
\label{AppendixISS}
This appendix uses many computations that are very similar to the {ones}
in Section~\ref{s2}, but in a simpler way. Thus, in order to avoid writing two times the same thing and to keep the proof relatively short, some steps might be quicker in this appendix.
Let $T_{1}\!>\!0$ and to be chosen later on. As $(H_{0}(0),V_{0}(0))$ satisfies \eqref{fluvial0}, there exists \hbox{$\nu_{a}\!>\!0$} such that for $\nu\in(0,\nu_{a})$, $F((H_{1}^{0},V_{1}^{0})^{T})$
has two distinct nonzero eigenvalues. Recall that $F$ is given by \eqref{defF} and that $\nu$ is the bound on \hbox{$\lVert H_{1}^{0}-H_{0}(0),V_{1}^{0}-V_{0}(0)\rVert_{H^{2}(0,L)}$}.
Besides, from \eqref{targetsteady2}, the function $(H_{0}(t,\cdot),V_{0}(t,\cdot))$ is the solution of a system of {ODEs} with an initial condition depending on a parameter $t$. Thus, as $\partial_{t}Q_{0}\in C^{2}([0,+\infty))$ and the slope $C$ satisfies $C\in C^{2}([0,L])$,
using \eqref{targetsteady0} and \cite[Chap.\,5, Th.\,3.1]{Hartman}, $(H_{0},V_{0})\in C^{3}([0,T_{1}]; C^{2}([0,L]))$ and there exists a constant $C$ depending only on $H_{\max}$, $\alpha$ {and
an upper bound of $\delta$,} such that,
{\begin{equation}
\label{boundH0V0}
\lVert\partial_{t}^{i}H_{0},\partial_{t}^{i}V_{0}\rVert_{C^{2}([0,L])}\leq C\sum\limits_{n=1}^{i}\left|\partial_{t}^{n}Q_{0}\right|,\quad\forall i\in[1,3],\ \forall t\in[0,T_{1}],
\end{equation} and in particular}
\begin{equation}
\lVert\partial_{t} H_{0},\partial_{t} V_{0}\rVert_{C^{2}([0,T_{1}];C^{2}([0,L]))}\leq C \lVert\partial_{t}Q_{0}\rVert_{C^{{2}}([0,+\infty))}.
\label{H0V0bound}
\end{equation}
Thus \cite[Th.\,2.1]{Wang} can still be used
{on $(H_{1}-H_{0})$}
and
there exist $\delta_{0}(T_{1})>0$ and $\nu_{0}(T_{1})\in(0,\nu_{a})$ such that, if $\nu\in(0,\nu_{0}(T_{1}))$ and $\delta\in(0,\delta_{0}(T_{1}))$, there exists a unique solution $(H_{1},V_{1})\in C^{0}([0,T_{1}];H^{2}(0,L))^{{2}}$ to the system \eqref{target}--\eqref{initialtarget}.
Besides $(H_{1},V_{1})$ satisfies an estimate as \eqref{estimate} but with $(H_{1},V_{1})$ instead of $(H,V)$ and $(H_{0},V_{0})$ instead of $(H_{1},V_{1})$. We denote by $C(T_{1})$ the associated constant.
Let us define $h_{1}:=H_{1}-H_{0}$ and $v_{1}:=V_{1}-V_{0}$.
We transform $(h_{1},v_{1})^{T}$ into $\boldsymbol{w}=(w_{1},w_{2})^{T}$ using the change of variables defined by
\eqref{change1}--\eqref{change2} with $H_{0}$ and $V_{0}$ instead of $H_{1}$ and $V_{1}$. Thus we obtain
\begin{equation}
\begin{split}
\partial_{t}\boldsymbol{w}+A_{0}(\boldsymbol{w},x)\partial_{x}\boldsymbol{w}&+B_{0}(\boldsymbol{w},x)+S_{0}\begin{pmatrix}\partial_{t}H_{0}\\ \partial_{t}V_{0}\end{pmatrix}=0,\\
w_{1}(t,0)&=\mathcal{H}_{1}(w_{2}(t,0),Q_{0}(t)-Q_{0}(0)),\\
w_{2}(t,L)&=\mathcal{H}_{2}(w_{2}(t,L)),
\end{split}
\label{sysw}
\end{equation}
where $A_{0}$, $B_{0}$ and $S_{0}$ have the same expression as $A$, $B$ and $S$ (given by\eqref{defA}, \eqref{defB}, \eqref{defS}) but with $(H_{0},V_{0})$ instead of $(H_{1}, V_{1})$. Similarly we define
\begin{equation*}
\lambda_{1}^{0}=V_{0}+\sqrt{gH_{0}},\quad\lambda_{2}^{0}=\sqrt{gH_{0}}-V_{0},
\end{equation*}
and $\phi^{0}$, defined as $\phi$ but with $(H_{0},V_{0})$ instead of $(H_{1}, V_{1})$.
Similarly as in Appendix~\ref{boundary},
\begin{equation*}
\mathcal{H}_{2}'(0)=-\lambda_{1}^{0}(L)/\lambda_{2}^{0}(L),\quad \mathcal{H}_{1}'(0)=-\lambda_{2}^{0}(0)/\lambda_{1}^{0}(0),
\end{equation*}
which is of the form \eqref{bound} with $v_{G}=0$ and $Z=0$.
Before going any further, note that we can perform the same computations as in Section 2 with no problem,
as the proof in Section \ref{s2} only used Proposition \ref{propISS} to get that $(H_{1},V_{1})$ exists for any time and that \eqref{fluvial} and Lemma \ref{lem3} hold, but we will see now that such claims are true for $H_{0}$ and $V_{0}$.
The existence of $(H_{0},V_{0})$ was already shown in Section~\ref{s1} and \eqref{fluvial0} is exactly \eqref{fluvial} with $(H_{0},V_{0})$ instead of $(H_{1},V_{1})$. Finally, \eqref{H0V0bound} is exactly the equivalent of Lemma \ref{lem3} for $(H_{0},V_{0})$.
We define now the Lyapunov function candidate $V:=V_{a}(\boldsymbol{w}(t,x),t)+V_{b}(\boldsymbol{w}(t,x),t)+V_{c}(\boldsymbol{w}(t,x),t)+V_{d}(\boldsymbol{w}(t,x),t)$, where $V_{a}$, $V_{b}$ and $V_{c}$ are defined in \eqref{Va}, \eqref{V2V3},
with $f_{1}$ and $f_{2}$ chosen as $f_{1}:=(\phi_{1}^{0})^{2}/(\lambda_{1}^{0}\eta)$ and $f_{2}:=(\phi_{2}^{0})^{2}\eta/(\lambda_{2}^{0})$,
where $\eta$ is a function such that there exists a constant $\varepsilon>0$ independent of $\boldsymbol{w}$ such that
\begin{equation*}
\begin{split}
\eta'&=\Bigl|\frac{\gamma_{2}^{0}}{\lambda_{1}^{0}}+\frac{\delta_{1}^{0}}{\lambda_{2}^{0}}\,\eta^{2}\Bigr|+\varepsilon,\quad \forall x\in[0,L],\\
\eta(0)&=\frac{\lambda_{2}^{0}(0)}{\lambda_{1}^{0}(0)}\,\phi^{0}(0)+\varepsilon.
\end{split}
\end{equation*}
Note that $\eta$ exists as, for any $t\in[0,+\infty)$, $(\phi(t,\cdot)^{0}\lambda_{2}^{0}(t\cdot)/\lambda_{1}^{0}(t\cdot))$ is a solution of
\begin{equation*}
\partial_{x} \chi=\Bigl|\frac{\gamma_{2}^{0}}{\lambda_{1}^{0}}+\frac{\delta_{1}^{0}}{\lambda_{2}^{0}}\,\chi^{2}\Bigr|, \quad\forall x\in[0,L],
\end{equation*}
this can be proved as in Lemma \ref{lem2}, and this case was actually shown in \cite{HS}.
Note that from \eqref{targetsteady0}, \eqref{targetsteady2} and \eqref{fluvial0}, $(H_{0})_{x}$ and $(V_{0})_{x}$ can be bounded by above and by below by constants that only depend on $H_{\max}$, $\alpha$ and an upper bound of $Q_{0}$ (which can also be expressed only with $H_{\max}$, $\alpha$ from \eqref{fluvial0}). Therefore, looking at their definition, the function $f_{1}$ and $f_{2}$ can also be bounded by above and below by constants that only depend on $H_{\max}$, $\alpha$ and $\varepsilon$. Thus there exist $c_{1}>0$ and $c_{2}>0$ depending only on $H_{\max}$ and $\alpha$, $\varepsilon$ and $\mu$ such that
\begin{equation}
c_{1}\lVert h_{1}(t,\cdot),v_{1}(t,\cdot)\rVert_{H^{2}(0,L)}^{2}\leq V(t)\leq c_{2}\lVert h_{1}(t,\cdot),v_{1}(t,\cdot)\rVert_{H^{2}(0,L)}^{2},\quad\forall t\in[0,T_{1}].
\label{equiv2}
\end{equation}
Consequently, by differentiating $V$ exactly as in \eqref{diffVa0}--\eqref{dVtot}, and from \eqref{sysw}, we~obtain that there exists $\mu>0$, $\nu_{1}\in(0,\nu_{0}(T_{1}))$ and ${\delta_{3}}>0$ such that, for any $\lVert h_{1}(0,\cdot),v_{1}(0,\cdot)\rVert_{H^{2}(0,L)}\leq \nu_{1}$,
and $\lVert\partial_{t}Q_{0}\rVert_{C^{2}([0,\infty))}\leq \delta$, where $\delta\in(0,\delta_{3})$,
\begin{equation*}
\begin{split}
\dot V\leq& -\mu V
+ \int_{0}^{L}2f_{1}w_{1}\biggl(S_{0}\begin{pmatrix}\partial_{t}H_{0}\\ \partial_{t}V_{0}\end{pmatrix}\biggr)_{1}+2f_{2}w_{2}\biggl(S_{0}\begin{pmatrix}\partial_{t}H_{0}\\ \partial_{t}V_{0}\end{pmatrix}\biggr)_{2}dx,\\
&+ \int_{0}^{L}2f_{1}\partial_{t}w_{1}\biggl(S_{0}\begin{pmatrix}\partial_{tt}^{2}H_{0}\\ \partial_{tt}^{2}V_{0}\end{pmatrix}\biggr)_{1}+2f_{2}\partial_{t}w_{2}\biggl(S_{0}\begin{pmatrix}\partial_{tt}^{2}H_{0}\\ \partial_{tt}^{2}V_{0}\end{pmatrix}\biggr)_{2}dx,\\
&+ \int_{0}^{L}2f_{1}\partial_{tt}^{2}w_{1}\biggl(S_{0}\begin{pmatrix}\partial_{ttt}^{3}H_{0}\\ \partial_{ttt}^{3}V_{0}\end{pmatrix}\biggr)_{1}+2f_{2}\partial_{tt}^{2}w_{2}\biggl(S_{0}\begin{pmatrix}\partial_{ttt}^{3}H_{0}\\ \partial_{ttt}^{3}V_{0}\end{pmatrix}\biggr)_{2}dx.
\end{split}
%\label{additional}
\end{equation*}
Thus, using Cauchy-Schwarz inequality, \eqref{equiv2}, {and \eqref{boundH0V0}} there exists $C_{1}>0$
depending only on $H_{\max}$, $\alpha$ and an upper bound of $\mu$ such that for any $t\in[0,T_{1}]$,
\begin{equation}
\label{H1V1decroissance}
\dot V(t)\leq -\mu V(t)
+C_{1}\left(|\partial_{t}Q_{0}(t)|+|\partial_{tt}^{2}Q_{0}(t)|+|\partial_{ttt}^{3}Q_{0}(t)|\right)V^{1/2}(t),
\end{equation}
and in particular
\begin{equation}
\label{H1V1decroissance1}
\dot V(t)\leq -\mu V(t)
+C_{1}\lVert\partial_{t}Q_{0}\rVert_{C^{2}([0,{t]})}V^{1/2}(t).
\end{equation}
Let us define $V_\mathrm{eq}:=(C_{1}\delta/\mu)^{2}$. From \eqref{H1V1decroissance1}, if $V(t)>2V_\mathrm{eq}$, then there exists a constant $k\!>\!0$ such that $\dot V(t)\!<\!-k V^{1/2}(t)$.
We now choose $\delta$ such that \hbox{${\sqrt{2}C_{1}\delta/(\mu \sqrt{c_{1}}})\!<\!\nu_{1}$}. Thus,
from \eqref{H1V1decroissance1} and as $c_{1}$, $c_{2}$, $C_{1}$ and $\mu$ do not depend on $T_{1}$, we can choose $T_{1}$ large enough such that
\begin{equation*}
V(T_{1})\leq 2V_\mathrm{eq}\leq c_{1}\nu_{1}^{2},
\end{equation*}
which implies that
\begin{equation*}
\lVert h_{1}(T_{1},\cdot),v_{1}(T_{1},\cdot)\rVert_{C^{2}(0,L)}\leq \nu_{1}
\end{equation*}
and therefore there exists a unique solution $(h_{1},v_{1})\!\in\! C^{0}([T_{1},2T_{1}],H^{2}(0,L))$, with initial condition $(h_{1}(T_{1},\cdot),v_{1}(T_{1},\cdot))$ (we use the same existence theorem (\cite[Th.\,2.1]{Wang})) and,
noting that $V(T_{1})\leq 2V_\mathrm{eq}$ implies $V(2T_{1})\leq 2V_\mathrm{eq}$,
this analysis still holds.
We~can do similarly for any $[nT_{1},(n+1)T_{1}]$ with $n\in\mathbb{N}$, thus, as
\[
(H_{0},V_{0})\in C^{0}([0,+\infty),H^{2}(0,L)),
\]
there exists a unique solution
\[
(H_{1},V_{1})\in C^{0}([0,+\infty),H^{2}(0,L))
\]
and \eqref{H1V1decroissance} holds for any $t\in[0,+\infty)$.
Therefore, denoting $g(t)=V(t)e^{\mu t}$,
we deduce from \eqref{H1V1decroissance} that
\begin{equation}
\label{ISS0}
g'(t)\leq C_{1}\left(|\partial_{t}Q_{0}(t)|+|\partial_{tt}^{2}Q_{0}(t)|+|\partial_{ttt}^{3}Q_{0}(t)|\right)e^{\sfrac{\mu t}{2}}\sqrt{g(t)}.
\end{equation}
Thus
\begin{equation*}
V^{1/2}(t)\leq V^{1/2}(0)e^{-\sfrac{\mu t}{2}}+\frac{C_{1}}{2} \biggl(\int_{0}^{t}\hspace*{-2mm}\bigl(|\partial_{t}Q_{0}(t)|+|\partial_{tt}^{2}Q_{0}(t)|+|\partial_{ttt}^{3}Q_{0}(t)|\bigr)e^{\sfrac{\mu s}{2}}ds\biggr)e^{-\sfrac{\mu t}{2}}.
\end{equation*}
This implies the ISS property
\begin{multline}
\label{eqISS}
\lVert h_{1}(t,\cdot),v_{1}(t,\cdot)\rVert_{H^{2}((0,L);\mathbb{R}^{2})}\leq \sqrt{\sfrac{c_{2}}{c_{1}}}\,\lVert h_{1}(0,\cdot),v_{1}(0,\cdot)\rVert_{H^{2}((0,L);\mathbb{R}^{2})}e^{-\sfrac{\mu t}{2}}\\
+\frac{C_{1}}{2\sqrt{c_{1}}}\biggl(\int_{0}^{t}\bigl(|\partial_{t}Q_{0}(t)|+|\partial_{tt}^{2}Q_{0}(t)|+|\partial_{ttt}^{3}Q_{0}(t)|\bigl)e^{\sfrac{\mu s}{2}}ds \biggr)e^{-\sfrac{\mu t}{2}}.
\end{multline}
This ends the proof of Proposition \ref{propISS}.
To extend this proof to the $H^{p}$ norm for $p>2$, note that using the same argument
\eqref{H0V0bound} holds with the $C^{p}([0,T_{1}];C^{3}([0,L]))$ norm in the left-hand side and the $C^{p}$ norm in the right-hand side.
We can can define $V_{3},\dots,V_{p}$ on $H^{p}(0,L)\times\mathbb{R}\times\mathbb{R}_{+}$ as in \eqref{V2V3} such that $V_{k}(\boldsymbol{w}(t,x),t)=V_{a}(\partial_{t}^{k}\boldsymbol{w}(t,x),t)$, for any $k\in[3,p]$.
Then \eqref{equiv2} holds with $V:=V_{a}+V_{b}+V_{c}+V_{3}+\cdots+ V_{p}$ and the $H^{p}$ norm, and the rest can done done identically.

\section{Proof of Theorem \ref{th2}}
\label{ISS2}
Theorem \ref{th2} result from the proof of Theorem \ref{th1}. Note that the boundary conditions \eqref{boundary2} can be written under the form \eqref{boundary1} with $(H_{0},V_{0})$ instead of $(H_{1},V_{1})$
where the only difference is that $Z$ satisfies now
\begin{equation}
\dot Z=H_{c}-H(t,L)+\frac{f(t)}{v_{G}k_{I}},
\label{Z2}
\end{equation}
where $f(t)=H_{c}\partial_{t}V_{0}(t,L)$.
The rest of the proof can be conducted as in Appendix~\ref{AppendixISS} for $(H_{1},V_{1})$,
with a priori two differences: $(H,V)$ satisfies the boundary conditions of the form \eqref{boundary1} and not of the form given in \eqref{target},
and $\dot Z$ satisfies \eqref{Z2} instead of~\eqref{Z}. However,
note that in Appendix \ref{AppendixISS} the only assumption used on the boundary conditions of the transformed system is that they are of the form \eqref{bound1},
which is still the case here.
Thus,
the only difference with Appendix \ref{AppendixISS} are
some addi\-tional terms when~$\dot Z$ is used, which is in the boundary terms in the derivative of the Lyapunov function. There exists therefore $\delta_{4}>0$ and $\nu_{2}>0$ such that{, for any $\lVert h_{1}(0,\cdot),v_{1}(0,\cdot)\rVert_{H^{2}(0,L)}\leq \nu_{2}$,
and $\lVert\partial_{t}Q_{0}\rVert_{C^{2}([0,\infty))}\leq \delta$, where $\delta\in(0,\delta_{4})$,}
\begin{multline*}
\dot V(t)\leq -\frac{\gamma}{2}\,V(t)
{+C_{1}\bigl|\partial_{t}Q_{0}(t)+\partial_{tt}^{2}Q_{0}(t)+\partial_{ttt}^{3}Q_{0}(t)\bigr|V^{1/2}}\\
+2qZf(t)+2q\dot Zf'(t)+2q\ddot Zf''(t),
\end{multline*}
where $C_{1}$ is a constant only depending on $H_{\max}$, $\alpha$, $\nu_{2}$ and $\delta_{4}$. Using Lemma \ref{lem3},
there exists a constant $C>0$ depending only on $H_{\max}$, $\alpha$, $\nu_{2}$ and $\delta_{4}$ such that
\begin{equation*}
\dot V\leq -\frac{\gamma}{2}\,V+CV^{1/2}\bigl|\partial_{t}Q_{0}(t)+\partial_{tt}^{2}Q_{0}(t)+\partial_{ttt}^{3}Q_{0}(t)\bigr|.
\end{equation*}
The same argument as in Appendix \ref{AppendixISS}, \eqref{ISS0}--\eqref{eqISS}, implies directly the ISS property~\eqref{ISSHV}.

\section{Proof of Lemma \ref{lem2}}
\label{sec:prooflem2}
In this appendix we prove Lemma \ref{lem2}. The proof is very similar to the {one} given in \cite{HS} in the special case where $(H_{1},V_{1})$ is a steady state. However, it happens that the proof actually does not need the relation $(H_{1}V_{1})_{x}=0$, which is no longer true when $(H_{1},V_{1})$ is not a steady-state. Let $\chi=(\lambda_{2}\phi/\lambda_{1})$, we have from \eqref{phi}:
\begin{equation}
\begin{split}
\partial_{x}\chi&=\frac{\phi}{\lambda_{1}^{2}}\left(\lambda_{1}\partial_{x}\lambda_{2}-\lambda_{2}\partial_{x}\lambda_{1}+\lambda_{2}\gamma_{1}+\lambda_{1}\delta_{2}\right)\\[-3pt]
&=\frac{\phi}{\lambda_{1}^{2}}\biggl((V_{1}+\sqrt{gH_{1}}) \Bigl(-V_{1x}+\frac{\sqrt{gH_{1}}}{2H_{1}}H_{1x}\Bigr)\\[-3pt]
&\hspace*{.5cm}-(-V_{1}+\sqrt{gH_{1}})\Bigl(V_{1x}+\frac{\sqrt{gH_{1}}}{2H_{1}}H_{1x}\Bigr)\\[-3pt]
&\hspace*{.5cm}+(\sqrt{gH_{1}}-V_{1}) \Bigl(\frac{3}{4}\sqrt{\sfrac{g}{H_{1}}}\,H_{1x}+\frac{3}{4}V_{1x}+\frac{kV}{H_{1}}-\frac{kV_{1}^{2}}{2H_{1}^{2}}\sqrt{\sfrac{H_{1}}{g}}\, \Bigr)\\[-3pt]
&\hspace*{.5cm}+(V_{1}+\sqrt{gH_{1}}) \Bigl(-\frac{3}{4}\sqrt{\sfrac{g}{H_{1}}}\,H_{1x}+\frac{3}{4}V_{1x}+\frac{2kV}{H_{1}}+\frac{kV_{1}^{2}}{2H_{1}^{2}}\sqrt{\sfrac{H_{1}}{g}}\, \Bigr)\biggr)\\[-3pt]
&=\frac{\phi}{\lambda_{1}^{2}}\biggl(\sqrt{gH_{1}} \Bigl(-2V_{1x}+\frac{3}{2}V_{1x}+\frac{2kV}{H_{1}} \Bigr)\\[-3pt]
&\hspace*{.5cm}-V_{1} \Bigl(\frac{3}{2}\sqrt{\sfrac{g}{H_{1}}}\,H_{1x}-\frac{kV_{1}^{2}}{H_{1}^{2}}\sqrt{\sfrac{H_{1}}{g}}\,-\sqrt{\sfrac{g}{H_{1}}}\,H_{1x} \Bigr)\biggr)\\[-3pt]
&=\frac{\phi}{\lambda_{1}^{2}} \Bigl(\frac{2kV}{H_{1}}\sqrt{gH_{1}}+\frac{kV_{1}^{2}}{H_{1}^{2}}\sqrt{\sfrac{H_{1}}{g}}\,V_{1}+\frac{1}{2}\sqrt{\sfrac{g}{H_{1}}}\,\partial_{t}H_{1} \Bigr).
\end{split}
\label{Le1}
\end{equation}
{On} the other hand:
\begin{equation}
\begin{split}
\Bigl(\frac{\phi\gamma_{2}}{\lambda_{1}}&+\frac{\delta_{1}}{\lambda_{2}\phi}\,\chi^{2}\Bigr)\\[-3pt]
&=\frac{\phi}{\lambda_{1}^{2}}\left(\lambda_{1}\gamma_{2}+\lambda_{2}\delta_{1}\right)\\[-3pt]
&=\frac{\phi}{\lambda_{1}^{2}} \Bigl(\frac{2kV}{H_{1}}\sqrt{gH_{1}}+\frac{kV_{1}^{2}}{H_{1}^{2}}\sqrt{\sfrac{H_{1}}{g}}\,V_{1}+V_{1}\sqrt{\sfrac{g}{H_{1}}}\,\frac{H_{1x}}{2}+V_{1x}\frac{\sqrt{gH_{1}}}{2} \Bigr)\\[-3pt]
&=\frac{\phi}{\lambda_{1}^{2}} \Bigl(\frac{2kV}{H_{1}}\sqrt{gH_{1}}+\frac{kV_{1}^{2}}{H_{1}^{2}}\sqrt{\sfrac{H_{1}}{g}}\,V_{1}-\frac{1}{2}\sqrt{\sfrac{g}{H_{1}}}\,\partial_{t}H_{1} \Bigr).
\end{split}
\label{Le2}
\end{equation}
Thus from \eqref{Le1} and \eqref{Le2}
\begin{equation}
\partial_{x}\chi= \Bigl(\frac{\phi\gamma_{2}}{\lambda_{1}}+\frac{\delta_{1}}{\lambda_{2}\phi}\,\chi^{2}+{\frac{\phi}{\lambda_{1}^{2}}}\sqrt{\sfrac{g}{H_{1}}}\,\partial_{t}H_{1} \Bigr),
\label{Leq}
\end{equation}
{and} there exists $\delta_{0}$ such that, if $\lVert \partial_{t}H_{1}\rVert_{L^{\infty}((0,+\infty)\times(0,L)}\leq \delta_{0}$,
\begin{equation*}
\frac{\phi}{\lambda_{1}^{2}} \Bigl(\frac{2kV_{1}}{H_{1}}\sqrt{gH_{1}}+\frac{kV_{1}^{2}}{H_{1}^{2}}\sqrt{\sfrac{H_{1}}{g}}\,V_{1}+\sqrt{\sfrac{g}{H_{1}}}\,\partial_{t}H_{1} \Bigr)>0,\quad\forall x\in[0,L],\ t\in[0,+\infty),
\end{equation*}
and, from \eqref{Le1} and \eqref{Leq},
\begin{equation*}
\partial_{x}\chi= \Bigl |\frac{\phi\gamma_{2}}{\lambda_{1}}+\frac{\delta_{1}}{\lambda_{2}\phi}\chi^{2}+{\frac{\phi}{\lambda_{1}^{2}}}\sqrt{\sfrac{g}{H_{1}}}\,\partial_{t}H_{1} \Bigr|,
\end{equation*}
this ends the proof of Lemma \ref{lem2}.

\section{Proof of Lemma \ref{lem3}}\label{app:prooflem3}
In this appendix we show that Lemma \ref{lem3} is a consequence of Proposition \ref{propISS} and Remark \ref{rmkHpISS}.
\begin{proof}
Indeed using Proposition \ref{propISS} and Remark \ref{rmkHpISS} with $p=3$, we have
\begin{multline*}
\|H_{1}(t,\cdot)-H_{0}(t,\cdot)\|_{H^{3}(0,L)}+\|V_{1}(t,\cdot)-V_{0}(t,\cdot)\|_{H^{3}(0,L)}\\
\leq \left(\|H_{1}^{0}-H^{*}\|_{H^{3}(0,L)}+\|V_{1}^{0}-V^{*}\|_{H^{3}(0,L)}\right)e^{-\sfrac{\mu t}{2}}\\
+c_{2} \frac{2}{\mu}(1-e^{-\sfrac{\mu t}{2}}) \|\partial_{t}Q_{0}\|_{C^{3}([0,+\infty))}.
\end{multline*}
Note that we chose $H_{1}^{0} = H^{*}$ and $V_{1}^{0} = V^{*}$ which means that
\begin{equation}\label{estimlemma1}
\|H_{1}(t,\cdot)-H_{0}(t,\cdot)\|_{H^{3}(0,L)}+\|V_{1}(t,\cdot)-V_{0}(t,\cdot)\|_{H^{3}(0,L)}\leq c_{2} \frac{2}{\mu}\|\partial_{t}Q_{0}\|_{C^{3}([0,+\infty))}.
\end{equation}
Note that $H_{1}-H_{1}^{0}$ is the solution of a quasilinear hyperbolic system and is small in~$H^{3}$ norm provided that $\partial_{t}Q_{0}$ is small in~$C^{3}$ norm. Therefore, there exists a constant~$C$ depending only on the parameters of the system and the bound $\nu$ such that
\begin{align}
\label{estimlemma2}
\|H_{1}&(t,\cdot)-H_{0}(t,\cdot)\|_{L^{2}(0,L)}+\|V_{1}(t,\cdot)-V_{0}(t,\cdot)\|_{L^{2}(0,L)}\\\notag
&+\|\partial_{t}H_{1}(t,\cdot)-\partial_{t}H_{0}(t,\cdot)\|_{L^{2}(0,L)}+\|\partial_{t}V_{1}(t,\cdot)-\partial_{t}V_{0}(t,\cdot)\|_{L^{2}(0,L)}\\\notag
&+\|\partial_{tx}^{2}H_{1}(t,\cdot)-\partial_{tx}^{2}H_{0}(t,\cdot)\|_{L^{2}(0,L)}+\|\partial_{tx}^{2}V_{1}(t,\cdot)-\partial_{tx}^{2}V_{0}(t,\cdot)\|_{L^{2}(0,L)}\\\notag
&+\|\partial_{tt}^{2}H_{1}(t,\cdot)-\partial_{tt}^{2}H_{0}(t,\cdot)\|_{L^{2}(0,L)}+\|\partial_{tt}^{2}V_{1}(t,\cdot)-\partial_{tt}^{2}V_{0}(t,\cdot)\|_{L^{2}(0,L)}\\\notag
&+\|\partial_{ttt}^{3}H_{1}(t,\cdot)-\partial_{ttt}^{3}H_{0}(t,\cdot)\|_{L^{2}(0,L)}+\|\partial_{ttt}^{3}V_{1}(t,\cdot)-\partial_{ttt}^{3}V_{0}(t,\cdot)\|_{L^{2}(0,L)}\\\notag
&+\|\partial_{ttx}^{3}H_{1}(t,\cdot)-\partial_{ttx}^{3}H_{0}(t,\cdot)\|_{L^{2}(0,L)}+\|\partial_{ttx}^{3}V_{1}(t,\cdot)-\partial_{ttx}^{3}V_{0}(t,\cdot)\|_{L^{2}(0,L)}
\\\notag
&\hspace*{2cm}\leq C\|H_{1}(t,\cdot)-H_{0}(t,\cdot)\|_{H^{3}(0,L)}+\|V_{1}(t,\cdot)-V_{0}(t,\cdot)\|_{H^{3}(0,L)}.
\end{align}
In what follows, the value of $C$ might change between lines but it still denotes a constant that only depends on the parameters of the system and the bound $\nu$. Besides, from Sobolev inequality, for $f\in H^{1}([0,L])$,
\begin{equation}\label{estimlemma3}
\|f\|_{C^{0}([0,L])} \leq C\left(\|f\|_{L^{2}([0,L])}+\|\partial_{x}f\|_{L^{2}([0,L])}\right).
\end{equation}
Combining \eqref{estimlemma1}, \eqref{estimlemma2} and \eqref{estimlemma3},
\begin{align*}
\max\bigl(&\|H_{1}(t,\cdot)-H_{0}(t,\cdot)\|_{C^{0}([0,L])},\|V_{1}(t,\cdot)-V_{0}(t,\cdot)\|_{C^{0}([0,L])}\bigr)\\
&+\max\bigl(\|\partial_{t}H_{1}(t,\cdot)-\partial_{t}H_{0}(t,\cdot)\|_{L^{\infty}(0,L)},\|\partial_{t}V_{1}(t,\cdot)-\partial_{t}V_{0}(t,\cdot)\|_{C^{0}([0,L])}\bigr)\\
&+
\max\bigl(\|\partial_{tt}^{2}H_{1}(t,\cdot)-\partial_{tt}^{2}H_{0}(t,\cdot)\|_{L^{\infty}(0,L)},\|\partial_{tt}^{{2}}V_{1}(t,\cdot)-\partial_{tt}^{{2}}V_{0}(t,\cdot)\|_{C^{0}([0,L])}\bigr)\\
&\hspace*{8cm}\leq Cc_{2} \frac{2}{\mu} \|\partial_{t}Q_{0}\|_{C^{3}([0,+\infty))}.
\end{align*}
Therefore, using the {reverse} triangular inequality and the fact that
\[
\max \bigl(\|H_{1}(t,\cdot)-H_{0}(t,\cdot)\|_{C^{0}([0,L])},\|V_{1}(t,\cdot)-V_{0}(t,\cdot)\|_{C^{0}([0,L])} \bigr)>0,
\]
we have
\begin{equation*}
\begin{split}
\max\bigl(&\|\partial_{t}H_{1}(t,\cdot)\|_{L^{\infty}(0,L)},\|\partial_{t}V_{1}(t,\cdot)\|_{L^{\infty}(0,L)}\bigr)\\&+
\max\bigl(\|\partial_{tt}^{2}H_{1}(t,\cdot)\|_{L^{\infty}(0,L)},\|\partial_{tt}^{{2}}V_{1}(t,\cdot)\|_{L^{\infty}(0,L)}\bigr)\\
\leq \max\bigl(&\|\partial_{t}H_{0}(t,\cdot)\|_{L^{\infty}(0,L)},\|\partial_{t}V_{0}(t,\cdot)\|_{L^{\infty}(0,L)}\bigr)\\
&+
\max\bigl(\|\partial_{tt}^{2}H_{0}(t,\cdot)\|_{L^{\infty}(0,L)},\|\partial_{tt}^{{2}}V_{0}(t,\cdot)\|_{L^{\infty}(0,L)}\bigr)+ Cc_{2} \frac{2}{\mu} \|\partial_{t}Q_{0}\|_{C^{3}([0,+\infty))}.
\end{split}
\end{equation*}
Recall that $(H_{0}(t,\cdot),V_{0}(t,\cdot))$ satisfies \eqref{boundH0V0},
{This implies}
\begin{multline*}
\max\bigl(\|\partial_{t}H_{1}(t,\cdot)\|_{L^{\infty}(0,L)},\|\partial_{t}V_{1}(t,\cdot)\|_{L^{\infty}(0,L)}\bigr)\\
+\max\bigl(\|\partial_{tt}^{2}H_{1}(t,\cdot)\|_{L^{\infty}(0,L)},\|\partial_{tt}V_{1}(t,\cdot)\|_{L^{\infty}(0,L)}\bigr)
\leq C\Bigl(1+c_{2} \frac{2}{\mu}\Bigr) \|\partial_{t}Q_{0}\|_{C^{3}([0,+\infty))}.
\end{multline*}
As this is true for any $t\in [0,+\infty)$ we have
\begin{equation*}
\|\partial_{t}H_{1},\partial_{t}V_{1}\|_{C^{1}([0,+\infty),C^{0}([0,L]))}\leq C\Bigl(1+c_{2} \frac{2}{\mu}\Bigr) \|\partial_{t}Q_{0}\|_{C^{3}([0,+\infty))}.
\end{equation*}
This ends the proof of Lemma \ref{lem3}.
\end{proof}

\section{Expression of \texorpdfstring{$V_{b}$ and $V_{c}$ for functions of $H^{2}(0,L)$}{V}}
\label{sec:VbVc}
Looking at \eqref{Va}, $V_{a}$ is indeed a function defined on $H^{2}(0,L)\times\mathbb{R}\times\mathbb{R}_{+}$. However, $V$ given by \eqref{def-V-tot} is a priori only defined for time-dependent functions $\boldsymbol{u}\in C^{2}([0,T],L^{2}(0,L))$. In fact, we can extend this do define $V$ on $H^{2}(0,L)\times\mathbb{R}\times \mathbb{R}_{+}$ by defining $V_{b}$ and $V_{c}$ as follows:
\begin{align}
V_{b}(\boldsymbol{U},t)&:=\int_{0}^{L}\hspace*{-2mm}f_{1}e^{-\mu x}(E(\boldsymbol{U}(x),x,t)\boldsymbol{I}(\boldsymbol{U},x,t))^{2}_{1}+f_{2}e^{\mu x}(E(\boldsymbol{U}(x),x,t)\boldsymbol{I}(\boldsymbol{U},x,t))^{2}_{2}dx\notag\\[-4pt]
&\hspace*{5.92cm}+q\frac{H_{1}(t,L)}{4g}(U_{1}(L)-U_{2}(L))^{2},\notag\\[-5pt]\label{V2V3}\\[-5pt]
V_{c}(\boldsymbol{U},t)&:=\int_{0}^{L}\hspace*{-2mm}f_{1}e^{-\mu x}(E(\boldsymbol{U}(x),x,t)\boldsymbol{J}(\boldsymbol{U},x,t))^{2}_{1}+f_{2}e^{\mu x}(E(\boldsymbol{U}(x),x,t)\boldsymbol{J}(\boldsymbol{U},x,t))^{2}_{2}dx\notag\\
&\hspace*{2cm}+q\Bigl(\sqrt{\sfrac{H_{1}(t,L)}{4g}}\,(I_{1}(t,L)-I_{2}(t,L))\notag\\
&\hspace*{4.3cm}+\frac{\partial_{t}H_{1}(t,L)}{4}\sqrt{\sfrac{1}{gH_{1}(t,L)}}\,(U_{1}(L)-U_{2}(L)) \Bigr)^{2},\notag
\end{align}
where
\begin{equation}
\begin{split}
\boldsymbol{I}(\boldsymbol{U},x,t)&:=A(\boldsymbol{U},x,t)\partial_{x}\boldsymbol{U}+B(\boldsymbol{U},x,t),\\
\boldsymbol{J}(\boldsymbol{U},x,t)&:=A(\boldsymbol{U},x,t)\partial_{x}(-\boldsymbol{I}(\boldsymbol{U},x,t))+(\partial_{t}A(\boldsymbol{U},x,t)\\
&\hspace*{1cm}+\partial_{\boldsymbol{U}}A(\boldsymbol{U},x,t)\cdot(-\boldsymbol{I}(\boldsymbol{U},x,t)))\partial_{x}\boldsymbol{U}+\partial_{t}B(\boldsymbol{U},x,t)\\
&\hspace*{1cm}+(\partial_{\boldsymbol{U}}B(\boldsymbol{U},x,t))({-\boldsymbol{I}(\boldsymbol{U},x,t)}),
\end{split}
\label{IJ}
\end{equation}
Observe that{, for a solution $\boldsymbol{u}$ of \eqref{sys1}, these quantities reduces to
\[
\boldsymbol{I}(\boldsymbol{U},x,t)=-\partial_{t}\boldsymbol{u}(t,x)\quad\text{and}\quad\boldsymbol{J}(\boldsymbol{U},x,t)=-\partial_{t}^{2}\boldsymbol{u}(t,x).
\]
Hence, using the expression of $Z$ given by \eqref{Z}, the expressions of $V_{b}(\boldsymbol{u}(t,\cdot),t)$ and $V_{c}(\boldsymbol{u}(t,\cdot),t)$ reduces to}
\begin{align*}
V_{b}(\boldsymbol{u}(t,\cdot),t)&:=\int_{0}^{L}\hspace*{-2mm}f_{1}(t,x)e^{-\mu x}(E\partial_{t}\boldsymbol{u})_{1}^{2}(t,x)\!+\!f_{2}(t,x)e^{\mu x}(E\partial_{t}\boldsymbol{u})_{2}^{2}(t,x)dx\!+\!q(\dot Z(t))^{2},\\
V_{c}(\boldsymbol{u}(t,\cdot),t)&:=\int_{0}^{L}\hspace*{-2mm}f_{1}(t,x)e^{-\mu x}(E\partial_{tt}^{2}\boldsymbol{u})_{1}^{2}(t,x)\!+\!f_{2}(t,x)e^{\mu x}(E\partial_{tt}^{2}\boldsymbol{u})_{2}^{2}(t,x)dx\!+\!q(\ddot Z(t))^{2},
%\label{V2V31}
\end{align*}
{which is exactly the definition \eqref{VbVc2} given earlier and} justifies the expression chosen for \eqref{V2V3} and \eqref{IJ}.

\section{Proof of Proposition \ref{propf1f2}}
\label{sec:propf1f2}
Let $T>{t_{0}\geq0}$ and $(\boldsymbol{u}^{0},Z^{0})\in H^{2}(0,L)\times\mathbb{R}$ satisfying the compatibility condition \eqref{compat} and such that
\begin{equation*}
\left(\lVert \boldsymbol{u}^{0}\rVert_{H^{2}(0,L)}+|Z^{0}|\right)<\nu,
%\label{boundu}
\end{equation*}
where $\nu$ is a constant to be chosen later on but such that $\nu<\min(\nu_{2},\nu({T-t_{0}}))$.
Let $(\boldsymbol{u},Z)\in C^{0}([{t_{0}},T],H^{3}(0,L))\times C^{2}([t_{0},T])$ be a solution with initial condition $\boldsymbol{u}^{0}$, $Z_{0}$.
Let $\delta>0$ to be chosen later on and assume that
\begin{equation*}
\max\bigl(\lVert\partial_{t}H_{1}\rVert_{C^{1}([t_{0},\infty);{C^{0}([0,L])})},\lVert\partial_{t}V_{1}\rVert_{C^{1}([t_{0},\infty);{C^{0}([0,L])})}\bigr)<\delta.
%\label{boundHVapp}
\end{equation*}
{As this is the only assumption on $H_{1}$ and $V_{1}$, we can assume from now on that $t_{0}=0$ without loss of generality.}

We start now by dealing with $V_{a}$. Differentiating $t\to V_{a}(t)$ with respect to time, using \eqref{sys1}, \eqref{definv} and integrating by parts, one has
\begin{align}
\label{diffVa0}
\dot V_{a}&= -2\int_{0}^{L}f_{1}(t,x)e^{-\mu x}(E\boldsymbol{u})_{1}\left[(EA(\boldsymbol{u},x,t)\partial_{x} \boldsymbol{u})_{1}+(EB)_{1}(\boldsymbol{u},x,t)\right]\\
&\hspace*{1cm}+f_{2}(t,x)e^{\mu x}(E\boldsymbol{u})_{2}\left[(EA(\boldsymbol{u},x,t)\partial_{x} \boldsymbol{u})_{2}+(EB)_{2}(\boldsymbol{u},x,t)\right]dx\notag\\
&\hspace*{1cm}+\int_{0}^{L}\partial_{t}(f_{1})e^{-\mu x}(E\boldsymbol{u})_{1}^{2}+\partial_{t}(f_{2})e^{\mu x}(E\boldsymbol{u})_{2}^{2}dx
\notag\\
&\hspace*{1cm}+2\int_{0}^{L}f_{1}e^{-\mu x}(E\boldsymbol{u})_{1}\bigl((\partial_{t}E+\partial_{\boldsymbol{u}}E\cdot\partial_{t}\boldsymbol{u})\boldsymbol{u}\bigr)_{1}\notag\\
&\hspace*{4cm}+f_{2}e^{\mu x}\bigl((\partial_{t}E+\partial_{\boldsymbol{u}}E\cdot\partial_{t}\boldsymbol{u})\boldsymbol{u}\bigr)_{2}dx\notag\\
&\hspace*{1cm}+2qZ(t)\dot Z(t)\notag
\end{align}
\begin{align*}
&=-2\int_{0}^{L}f_{1}(t,x)e^{-\mu x}(E\boldsymbol{u})_{1}\left[D_{1}(\boldsymbol{u},x,t)\left(\partial_{x} (E\boldsymbol{u})-(\partial_{x}E+\partial_{\boldsymbol{U}}E\cdot\partial_{x}\boldsymbol{u})\boldsymbol{u}\right)_{1}\right]\\
&\hspace*{1cm}+f_{2}(t,x)e^{\mu x}(E\boldsymbol{u})_{2}\left[D_{2}(\boldsymbol{u},x,t)\left(\partial_{x} (E\boldsymbol{u})-(\partial_{x}E+\partial_{\boldsymbol{U}}E\cdot\partial_{x}\boldsymbol{u})\boldsymbol{u}\right)_{2}\right]dx\\
&\hspace*{1cm}+\int_{0}^{L}\partial_{t}(f_{1})e^{-\mu x}(E\boldsymbol{u})_{1}^{2}+\partial_{t}(f_{2})e^{\mu x}(E\boldsymbol{u})_{2}^{2}dx\\
&\hspace*{1cm}-2\int_{0}^{L}f_{1}e^{-\mu x}(E\boldsymbol{u})_{1}(EB)_{1}(\boldsymbol{u},x,t)+f_{2}e^{\mu x}(E\boldsymbol{u})_{2}(EB)_{2}(\boldsymbol{u},x,t)dx\\
&\hspace*{1cm}+2\int_{0}^{L}f_{1}e^{-\mu x}(E\boldsymbol{u})_{1}\bigl((\partial_{t}E+\partial_{\boldsymbol{u}}E\cdot\partial_{t}\boldsymbol{u})\boldsymbol{u}\bigl)_{1}\\
&\hspace*{4cm}+f_{2}e^{\mu x}\bigl((\partial_{t}E+\partial_{\boldsymbol{u}}E\cdot\partial_{t}\boldsymbol{u})\boldsymbol{u}\bigl)_{2}dx\\
&\hspace*{1cm}+2qZ(t)\dot Z(t),
\end{align*}
\begin{align}
\label{diffVa2}
\dot V_{a}={}& -\left[f_{1}e^{-\mu x}D_{1}(E\boldsymbol{u})_{1}^{2}+D_{2}f_{2}e^{\mu x}(E\boldsymbol{u})_{2}^{2}\right]_{0}^{L}\\ \notag
&-\int_{0}^{L}(E\boldsymbol{u})_{1}e^{-\mu x}\Bigl(\bigl(-\partial_{x}(D_{1}f_{1})-f_{1}\partial_{u}(D_{1})\cdot\partial_{x}\boldsymbol{u} \bigr)(E\boldsymbol{u})_{1}\\ \notag
&\hspace*{5cm}-2f_{1}D_{1} \bigl((\partial_{x}E+\partial_{\boldsymbol{U}}E\cdot\partial_{x}\boldsymbol{u})\boldsymbol{u} \bigl)_{1}\Bigr)\\ \notag
&+(E\boldsymbol{u})_{2}e^{\mu x}\Bigl(\bigl(-\partial_{x}(D_{2}f_{2})-f_{2}\partial_{u}(D_{2})\cdot\partial_{x}\boldsymbol{u})(E\boldsymbol{u} \bigr)_{2}\\ \notag
&\hspace*{5cm}-2f_{2}D_{2} \bigl((\partial_{x}E+\partial_{\boldsymbol{U}}E\cdot\partial_{x}\boldsymbol{u})\boldsymbol{u} \bigr)_{2}\Bigr)dx\\ \notag
&+\int_{0}^{L}\partial_{t}(f_{1})e^{-\mu x}(E\boldsymbol{u})_{1}^{2}+\partial_{t}(f_{2})e^{\mu x}(E\boldsymbol{u})_{2}^{2}dx\\ \notag
&-2\int_{0}^{L}f_{1}e^{-\mu x}(E\boldsymbol{u})_{1}(EB)_{1}(\boldsymbol{u},x,t)+f_{2}e^{\mu x}(E\boldsymbol{u})_{2}(EB)_{2}(\boldsymbol{u},x,t)\\ \notag
&+2\int_{0}^{L}f_{1}e^{-\mu x}(E\boldsymbol{u})_{1} \bigl((\partial_{t}E+\partial_{\boldsymbol{u}}E\cdot\partial_{t}\boldsymbol{u})\boldsymbol{u} \bigr)_{1}\\ \notag
&\hspace*{5cm}+f_{2}e^{\mu x} \bigl((\partial_{t}E+\partial_{\boldsymbol{u}}E\cdot\partial_{t}\boldsymbol{u})\boldsymbol{u} \bigr)_{2}dx\\ \notag
&-\mu\int_{0}^{L}D_{1}f_{1}e^{-\mu x}(E\boldsymbol{u})_{1}^{2}-D_{2}f_{2}e^{\mu x}(E\boldsymbol{u})_{2}^{2}dx
+2qZ(t)\dot Z(t).
\end{align}
In order to simplify this expression, observe that from \eqref{defA}, \eqref{definv} and \eqref{Lambda}, $D_{1}(\boldsymbol{0},x,t)=\lambda_{1}(t,x)$ and $D_{2}(\boldsymbol{0},x,t)=-\lambda_{2}(t,x)$. Recall that $H_{1}$ is bounded by $H_{\max}$ from \eqref{Hinfty} and that $gH_{1}-V_{1}^{2}$ is bounded by below by $\alpha$ from \eqref{fluvial}. Using this, the fact that
$D$ is $C^{1}$ in $\boldsymbol{u}$,
and using also \eqref{ineq} and \eqref{boundHV}, there exists $C>0$ depending only on $H_{\max}$ and $\alpha$, $\nu$ and $\delta$ such that
\begin{align*}
\lVert D_{i}-\sgn(D_{i})\lambda_{i}\rVert_{\infty}&\leq \lVert C\boldsymbol{u}\rVert_{\infty},
\\
\lVert\partial_{x}D_{i}+\partial_{\boldsymbol{u}}D_{i}.\partial_{x}\boldsymbol{u}-\sgn(D_{i})\partial_{x}\lambda_{i}\rVert_{\infty}&\leq C(\lVert \partial_{x}\boldsymbol{u}\rVert_{\infty}+\lVert \boldsymbol{u}\rVert_{\infty}),\quad i\in\{1,2\},
\end{align*}
and
\begin{align*}
\lVert\partial_{x}E\rVert_{\infty}&\leq C(\lVert \boldsymbol{u}\rVert_{\infty}),\\
\lVert\partial_{t}E+\partial_{\boldsymbol{u}}E\partial_{t}\boldsymbol{u}\rVert_{\infty}&\leq C(\lVert \boldsymbol{u}\rVert_{\infty}+\lVert \partial_{t} \boldsymbol{u}\rVert_{\infty}).
\end{align*}
Thus, using this together with \eqref{diffVa2}
\begin{equation*}
\begin{split}
\dot V_{a}\leq& -\left[f_{1}e^{-\mu x}D_{1}(E\boldsymbol{u})_{1}^{2}+D_{2}f_{2}e^{\mu x}(E\boldsymbol{u})_{2}^{2}\right]_{0}^{L}\\
&-\int_{0}^{L}(E\boldsymbol{u})_{1}^{2}e^{-\mu x}\bigl(-\partial_{x}(\lambda_{1}f_{1})-\partial_{t}(f_{1})\bigr)+(E\boldsymbol{u})_{2}^{2}e^{\mu x}\bigl(\partial_{x}(\lambda_{2}f_{2})-\partial_{t}(f_{2})\bigr)dx\\
&-2\int_{0}^{L}f_{1}e^{-\mu x}(E\boldsymbol{u})_{1}(EB)_{1}(\boldsymbol{u},x,t)+f_{2}e^{\mu x}(E\boldsymbol{u})_{2}(EB)_{2}(\boldsymbol{u},x,t)dx\\
&-\mu\int_{0}^{L}\lambda_{1}f_{1}e^{-\mu x}(E\boldsymbol{u})_{1}^{2}+\lambda_{2}f_{2}e^{\mu x}(E\boldsymbol{u})_{2}^{2}dx
+2qZ(t)\dot Z(t)\\
&+C\left(\lVert \boldsymbol{u}\rVert_{\infty}+\lVert \partial_{x} \boldsymbol{u}\rVert_{\infty}\right)\int_{0}^{L}(E\boldsymbol{u})_{1}^{2}+(E\boldsymbol{u})_{2}^{2} dx\\
&+C\left(\lVert \boldsymbol{u}\rVert_{\infty}+\lVert \partial_{x} \boldsymbol{u}\rVert_{\infty}\right)^{2}\int_{0}^{L}|(E\boldsymbol{u})_{1}|+|(E\boldsymbol{u})_{2}|dx,
\end{split}
\end{equation*}
where $C$ is a constant that may change between lines but
only depends on $\nu$, an upper bound of $\delta$ (for instance $\delta_{0}$), $\mu$, $H_{\max}$ and $\alpha$. Note that $C$ is continuous in $\mu\in[0,\infty)$, thus it can be made independent of $\mu$ by imposing an upper bound on $\mu$, for instance $\mu\in(0,1]$.
Finally, from the second equation of \eqref{definv}, and the fact that $E$ is $C^{1}$ in~$\boldsymbol{u}$,
there exists a continuous function {$\mathcal{E}_{1}$} defined on $\mathcal{B}_{\nu_{2}}\times[0,L]\times[0,T]$ such that, for any vector $\boldsymbol{v}\in\mathbb{R}^{2}$ and any $(t,x)\in[0,T]\times[0,L]$,
\begin{equation}
E(\boldsymbol{u}(t,x),x,t)\boldsymbol{v}-\boldsymbol{v}=(\boldsymbol{u}(t,x)\cdot\mathcal{E}_{1}(\boldsymbol{u}(t,x),x,t))\boldsymbol{v}.
\label{recoverv}
\end{equation}
As $E(\boldsymbol{u}(t,x),x,t)$ is a $C^{\infty}$ function of the coefficients of $A$, $\mathcal{E}_{1}$ is bounded on $\mathcal{B}_{\nu_{2}}\times[0,L]\times[0,T]$ by a bound that only depends on $\nu_{2}$, $H_{\max}$ and $\alpha$.
Thus there exists a constant $\bar C$ depending only on $\nu_{2}$, $H_{\max}$ and $\alpha$ such that
\begin{equation}
\frac{1}{\bar C}\,\lVert\boldsymbol{v}\rVert_{{L^{2}((0,L);\mathbb{R}^{2})}}\leq \lVert E\boldsymbol{v}\rVert_{{L^{2}((0,L);\mathbb{R}^{2})}}\leq \bar C\, \lVert\boldsymbol{v}\rVert_{{L^{2}((0,L);\mathbb{R}^{2})}}.
\label{equivinvu}
\end{equation}
Thus, using this together with the fact that $D_{1}$ and $D_{2}$ are $C^{1}$ with $\boldsymbol{u}$, \eqref{Va}, and Young's inequality and then Cauchy-Schwarz inequality on the last integral term,
\begin{equation*}
\begin{split}
\dot V_{a}\leq& -\left[f_{1}e^{-\mu x}\lambda_{1}(E\boldsymbol{u})_{1}^{2}-\lambda_{2}f_{2}e^{\mu x}(E\boldsymbol{u})_{2}^{2}\right]_{0}^{L}\\
&-\int_{0}^{L}(E\boldsymbol{u})_{1}^{2}e^{-\mu x}\bigl(-\partial_{x}(\lambda_{1}f_{1})-\partial_{t}(f_{1})\bigr)+(E\boldsymbol{u})_{2}^{2}e^{\mu x}\bigl(\partial_{x}(\lambda_{2}f_{2})-\partial_{t}(f_{2})\bigr)dx\\
&-2\int_{0}^{L}f_{1}e^{-\mu x}(E\boldsymbol{u})_{1}(EB)_{1}(\boldsymbol{u},x,t)+f_{2}e^{\mu x}(E\boldsymbol{u})_{2}(EB)_{2}(\boldsymbol{u},x,t)dx\\
&-\mu\min\limits_{x\in[0,L]}(\lambda_{1},\lambda_{2})V_{a}{+\mu\min\limits_{x\in[0,L]}(\lambda_{1},\lambda_{2})qZ^{2}(t)}
+2qZ(t)\dot Z(t)\\
&+C\left(\lVert \boldsymbol{u}\rVert_{\infty}+\lVert \partial_{x} \boldsymbol{u}\rVert_{\infty}\right)\lVert\boldsymbol{u} \rVert_{L^{2}(0,L)}^{2}\\
&+C\left(\lVert \boldsymbol{u}\rVert_{\infty}+\lVert \partial_{x} \boldsymbol{u}\rVert_{\infty}\right)^{3}+C\lVert \boldsymbol{u}\rVert_{\infty}(|\boldsymbol{u}(t,0)^{2}|+|\boldsymbol{u}(t,L)^{2}|).
\end{split}
\end{equation*}
Now, as $E$ and $B$ are $C^{2}$ with $\boldsymbol{u}$ and continuous with $x$ and $t$, and as $B(\boldsymbol{0},x,t)=0$, there exists a continuous function {$\mathcal{E}_{2}\in C^{0}(\mathcal{B}_{\nu_{2}}\times[0,T]\times[0,L];\mathbb{R}^{2\times 2\times 2})$} such that for any $(t,x)\in[0,T]\times[0,L]$
\begin{equation}
(EB)(\boldsymbol{u}(t,x),x,t)=\partial_{\boldsymbol{u}}(EB)(\boldsymbol{0},x,t)\cdot\boldsymbol{u}(t,x)+(\mathcal{E}_{2}(\boldsymbol{u},x,t)\cdot\boldsymbol{u}(t,x))\boldsymbol{u}(t,x).
\label{Bprop}
\end{equation}
Note that from \eqref{defB}, $\mathcal{E}_{2}$ is bounded on $\mathcal{B}_{\nu_{2}}\times[0,L]\times[0,T]$ by a constant that only depends on $\nu_{2}$, $\delta$, $H_{\max}$ and $\alpha$. From \eqref{defB} and \eqref{definv}
$\partial_{\boldsymbol{u}}(EB)(\boldsymbol{0},x,t)=\partial_{\boldsymbol{u}}B(\boldsymbol{0},x,t)$. Besides, from \eqref{definv}, $E$ is invertible and $C^{1}$, thus an inequality similar to
\eqref{equivinvu}
holds for $E^{-1}$, and $\boldsymbol{u}=E^{-1}(E\boldsymbol{u})$.
Therefore, using \eqref{Bprop} together with \eqref{recoverv}, {the fact that~$\mathcal{E}_{1}$ and $\mathcal{E}_{2}$ are bounded,} and the expression of $\partial_{\boldsymbol{u}}B(\boldsymbol{0},x,t)$ given in {\eqref{l1l2}}--\eqref{gamma}, one has
\begin{align*}
\dot V_{a}\leq& -\left[f_{1}e^{-\mu x}\lambda_{1}u_{1}^{2}-\lambda_{2}f_{2}e^{\mu x}u_{2}^{2}\right]_{0}^{L}\\
&-\int_{0}^{L}(E\boldsymbol{u})_{1}^{2}e^{-\mu x}\bigl(-\partial_{x}(\lambda_{1}f_{1})-\partial_{t}(f_{1})\bigr)+(E\boldsymbol{u})^{2}e^{\mu x}(\partial_{x}\bigl(\lambda_{2}f_{2})-\partial_{t}(f_{2})\bigr)dx\\
&-2\int_{0}^{L}\hspace*{-2mm}f_{1}e^{-\mu x}\gamma_{1}(E\boldsymbol{u})_{1}^{2}+f_{2}e^{\mu x}\delta_{2}(E\boldsymbol{u})_{2}^{2}+\left(\gamma_{2}f_{1}e^{-\mu x}+\delta_{1}f_{2}e^{\mu x}\right)(E\boldsymbol{u})_{1}(E\boldsymbol{u})_{2}dx\\
&-\mu\min\limits_{x\in[0,L]}(\lambda_{1},\lambda_{2})V_{a}{+\mu\min\limits_{x\in[0,L]}(\lambda_{1},\lambda_{2})qZ^{2}(t)}
+2qZ(t)\dot Z(t)\\
&+C\left(\lVert \boldsymbol{u}\rVert_{\infty}+\lVert \partial_{x} \boldsymbol{u}\rVert_{\infty}\right)\lVert\boldsymbol{u} \rVert_{L^{2}(0,L)}^{2}\\
&+C\left(\lVert \boldsymbol{u}\rVert_{\infty}+\lVert \partial_{x} \boldsymbol{u}\rVert_{\infty}\right)^{3}+C\lVert \boldsymbol{u}\rVert_{\infty}(|\boldsymbol{u}(t,0)^{2}|+|\boldsymbol{u}(t,L)^{2}|).
\end{align*}
As $\mathcal{D}_{1}$ and $\mathcal{D}_{2}$ are of class $C^{2}$,
denoting for simplicity
\[
k_{2}:=\partial_{1} \mathcal{D}_{1}(0,t),\quad k_{1}:=\partial_{1}\mathcal{D}_{2}(0,0,t),\quad
k_{3}:=-\partial_{2}\mathcal{D}_{2}(0,0,t),
\]
and using \eqref{bound}
\begin{multline*}
%\label{diffVa}
\dot V_{a}\leq
-\mu\min\limits_{x\in[0,L]}(\lambda_{1},\lambda_{2})V_{a}+\left[f_{1}\lambda_{1}k_{2}^{2}-\lambda_{2}f_{2}\right]u_{2}^{2}(t,0)\\
-I_{1}(u_{1}(t,L),Z(t))
-\int_{0}^{L}I_{2}((E\boldsymbol{u})_{1},(E\boldsymbol{u})_{2})dx+C(\lVert \boldsymbol{u}\rVert_{\infty}+\lVert \partial_{x} \boldsymbol{u}\rVert_{\infty})\\{}\cdot\Bigl(\lVert\boldsymbol{u} \rVert_{L^{2}(0,L)}^{2}+\left(\lVert \boldsymbol{u}\rVert_{\infty}+\lVert \partial_{x} \boldsymbol{u}\rVert_{\infty}\right)^{2}+(|\boldsymbol{u}(t,0)^{2}|+|\boldsymbol{u}(t,L)^{2}|)\Bigr),
\end{multline*}
where $I_{1}$ and $I_{2}$ denote the following quadratic forms
\begin{multline*}
I_{1}(x,y)=\bigl(\lambda_{1}f_{1}(L)e^{-\mu L}-\lambda_{2}f_{2}(L)e^{\mu L}k_{1}^{2}\bigr)x^{2}\\
+{\Bigl(q\sqrt{\sfrac{H_{1}}{g}}\,k_{3}-\lambda_{2}f_{2}(L)e^{\mu L}k_{3}^{2}-\mu\min\limits_{x\in[0,L]}(\lambda_{1},\lambda_{2})q\Bigr)}y^{2}\\
\shoveright{+\bigl(2\lambda_{2}f_{2}(L)e^{\mu L}k_{3}k_{1}-{q\sqrt{\sfrac{H_{1}}{g}}}\,(k_{1}-1)\bigr)xy,}\\
\shoveleft{I_{2}(x,y)=\bigl((-\lambda_{1}f_{1})_{x}+2f_{1}\gamma_{1}(t,x)-\partial_{t}f_{1}\bigr)e^{-\mu x}x^{2}}\\
+\bigl((\lambda_{2}f_{2})_{x}+2f_{2}\delta_{2}(t,x)-\partial_{t}f_{2}\bigr)e^{\mu x}y^{2}\\
+2\left(\gamma_{2}f_{1}e^{-\mu x}+\delta_{1}f_{2}e^{\mu x}\right)xy.
%\label{I1I2-app}
\end{multline*}
We can perform similarly with $V_{b}$ and $V_{c}$, to do this observe that $\partial_{t}\boldsymbol{u}$ and $\partial_{tt}^{2}\boldsymbol{u}$ are respectively solutions of
\begin{multline}
\label{V2}
\partial_{t}(\partial_{t}\boldsymbol{u})+A(\boldsymbol{u},x,t)\partial_{x}(\partial_{t}\boldsymbol{u})+(\partial_{\boldsymbol{u}}B(\boldsymbol{u},x,t))(\partial_{t}\boldsymbol{u})
\\
+(\partial_{t}A(\boldsymbol{u},x,t)+\partial_{\boldsymbol{u}}A(\boldsymbol{u},x,t).\partial_{t}\boldsymbol{u})\partial_{x}\boldsymbol{u}
+\partial_{t}B\left(\boldsymbol{u},x,t\right)=0,
\end{multline}
and
\begin{align}
\label{V3}
\partial_{t}(\partial_{tt}^{2}\boldsymbol{u})+A(\boldsymbol{u},x,t)\partial_{x}&(\partial_{tt}^{2}\boldsymbol{u})
+(\partial_{\boldsymbol{u}}A(\boldsymbol{u},x)\cdot\partial_{tt}^{2}\boldsymbol{u})\partial_{x}\boldsymbol{u}
+(\partial_{\boldsymbol{u}}B(\boldsymbol{u},x))(\partial_{tt}^{2}\boldsymbol{u})\\
\notag&+2\partial_{\boldsymbol{u}}(\partial_{t}A(\boldsymbol{u},x,t))\cdot\partial_{t}\boldsymbol{u})\partial_{x}\boldsymbol{u}+\partial_{tt}^{2}A(\boldsymbol{u},x,t)\partial_{x}\boldsymbol{u}\\
\notag&+2\partial_{\boldsymbol{u}}A(\boldsymbol{u},x)\cdot\partial_{t}\boldsymbol{u}\partial_{x}(\partial_{t}\boldsymbol{u})+\partial_{t}A(\boldsymbol{u},x,t)\partial_{x}(\partial_{t}\boldsymbol{u})\\
\notag&+((\partial^{2}_{\boldsymbol{u}}A(\boldsymbol{u},x)\cdot\partial_{t}\boldsymbol{u})\cdot\partial_{t}\boldsymbol{u})\partial_{x}\boldsymbol{u}+\partial_{tt}^{2}B(\boldsymbol{u},x)\\
\notag&+\partial_{\boldsymbol{u}}(\partial_{t}B(\boldsymbol{u},x))\cdot\partial_{t}\boldsymbol{u}
+(\partial^{2}_{\boldsymbol{u}}B(\boldsymbol{u},x)\cdot\partial_{t}\boldsymbol{u})(\partial_{t}\boldsymbol{u})=0,
\end{align}
which are very similar to \eqref{sys1}, as they only differ by quadratic perturbations or terms involving a time derivative of $(H_{1},V_{1})$.
We get then
\begin{align*}
\dot V=\dot V_{a}&+\dot V_{b}+\dot V_{c}\leq-\mu\min\limits_{x\in[0,L]}(\lambda_{1},\lambda_{2})V\\
&+\left[f_{1}\lambda_{1}k_{2}^{2}-\lambda_{2}f_{2}\right]\left(u_{2}^{2}(t,0)+(\partial_{t}u_{2}(t,0))^{2}+(\partial_{tt}^{2}u_{2}(t,0))^{2}\right)\\
&-I_{1}(u_{1}(t,L),Z)-I_{1}(\partial_{t}u_{1}(t,L),\dot Z)-I_{1}(\partial_{tt}^{2}u_{1}(t,L),\ddot Z)\\
&-\int_{0}^{L}\hspace*{-2mm}I_{2}((E\boldsymbol{u})_{1},(E\boldsymbol{u})_{2})+I_{2}((E\partial_{t}\boldsymbol{u})_{1},(E\partial_{t}\boldsymbol{u})_{2})+I_{2}((E\partial_{tt}^{2}\boldsymbol{u})_{1},(E\partial_{tt}^{2}\boldsymbol{u})_{2})dx\\
&+C(\lVert \boldsymbol{u}\rVert_{\infty}+\lVert \partial_{x} \boldsymbol{u}\rVert_{\infty})\Bigl(\lVert\boldsymbol{u} \rVert_{L^{2}(0,L)}^{2}+\lVert\partial_{t}\boldsymbol{u} \rVert_{L^{2}(0,L)}^{2}
+\lVert\partial_{tt}^{2}\boldsymbol{u} \rVert_{L^{2}(0,L)}^{2}\\
&\hspace*{1cm}+\left(\lVert \boldsymbol{u}\rVert_{\infty}+\lVert \partial_{x} \boldsymbol{u}\rVert_{\infty}\right)^{{2}}+|u_{2}(t,0)^{2}|+(|u_{1}(t,L)|+|Z|)^{2}+|\partial_{t}u_{2}(t,0)^{2}|\\
&\hspace*{2.3cm}+(|\partial_{t}u_{1}(t,L)|+|\dot Z|)^{2}+|\partial_{tt}^{2}u_{2}(t,0)^{2}|+(|\partial_{tt}^{2}u_{1}(t,L)|+|\ddot Z|)^{2}\Bigr)\\
&+C\delta\Bigl(|u_{2}(t,0)|^{2}+(|u_{1}(t,L)|+|Z|)^{2}+|\partial_{t}u_{2}(t,0)|^{2}\\
&\hspace*{5cm}+(|\partial_{t}u_{1}(t,L)|+|\dot Z|)^{2}\Bigr)+C\delta V.
%\label{dVtotapp}
\end{align*}

The two last terms come from the successive differentiations of the boundary conditions \eqref{bound}, together with \eqref{boundHV}, or from the terms in \eqref{V2}--\eqref{V3} involving a time derivative of $A$ or $B$. This ends the proof of Proposition \ref{propf1f2}.

